\documentclass[
]{article}
\usepackage{amsmath,amssymb}
\usepackage[letterpaper, margin=1in]{geometry}
\usepackage{iftex}
\ifPDFTeX
  \usepackage[T1]{fontenc}
  \usepackage[utf8]{inputenc}
  \usepackage{textcomp} % provide euro and other symbols
\else % if luatex or xetex
  \usepackage{unicode-math} % this also loads fontspec
  \defaultfontfeatures{Scale=MatchLowercase}
  \defaultfontfeatures[\rmfamily]{Ligatures=TeX,Scale=1}
\fi
\usepackage{lmodern}
\ifPDFTeX\else
  % xetex/luatex font selection
\fi
% Use upquote if available, for straight quotes in verbatim environments
\IfFileExists{upquote.sty}{\usepackage{upquote}}{}
\IfFileExists{microtype.sty}{% use microtype if available
  \usepackage[]{microtype}
  \UseMicrotypeSet[protrusion]{basicmath} % disable protrusion for tt fonts
}{}
\makeatletter
\@ifundefined{KOMAClassName}{% if non-KOMA class
  \IfFileExists{parskip.sty}{%
    \usepackage{parskip}
  }{% else
    \setlength{\parindent}{0pt}
    \setlength{\parskip}{6pt plus 2pt minus 1pt}}
}{% if KOMA class
  \KOMAoptions{parskip=half}}
\makeatother
\usepackage{xcolor}
\usepackage{color}
\usepackage{fancyvrb}
\newcommand{\VerbBar}{|}
\newcommand{\VERB}{\Verb[commandchars=\\\{\}]}
\DefineVerbatimEnvironment{Highlighting}{Verbatim}{commandchars=\\\{\}}
% Add ',fontsize=\small' for more characters per line
\newenvironment{Shaded}{}{}
\newcommand{\AlertTok}[1]{\textcolor[rgb]{1.00,0.00,0.00}{\textbf{#1}}}
\newcommand{\AnnotationTok}[1]{\textcolor[rgb]{0.38,0.63,0.69}{\textbf{\textit{#1}}}}
\newcommand{\AttributeTok}[1]{\textcolor[rgb]{0.49,0.56,0.16}{#1}}
\newcommand{\BaseNTok}[1]{\textcolor[rgb]{0.25,0.63,0.44}{#1}}
\newcommand{\BuiltInTok}[1]{\textcolor[rgb]{0.00,0.50,0.00}{#1}}
\newcommand{\CharTok}[1]{\textcolor[rgb]{0.25,0.44,0.63}{#1}}
\newcommand{\CommentTok}[1]{\textcolor[rgb]{0.38,0.63,0.69}{\textit{#1}}}
\newcommand{\CommentVarTok}[1]{\textcolor[rgb]{0.38,0.63,0.69}{\textbf{\textit{#1}}}}
\newcommand{\ConstantTok}[1]{\textcolor[rgb]{0.53,0.00,0.00}{#1}}
\newcommand{\ControlFlowTok}[1]{\textcolor[rgb]{0.00,0.44,0.13}{\textbf{#1}}}
\newcommand{\DataTypeTok}[1]{\textcolor[rgb]{0.56,0.13,0.00}{#1}}
\newcommand{\DecValTok}[1]{\textcolor[rgb]{0.25,0.63,0.44}{#1}}
\newcommand{\DocumentationTok}[1]{\textcolor[rgb]{0.73,0.13,0.13}{\textit{#1}}}
\newcommand{\ErrorTok}[1]{\textcolor[rgb]{1.00,0.00,0.00}{\textbf{#1}}}
\newcommand{\ExtensionTok}[1]{#1}
\newcommand{\FloatTok}[1]{\textcolor[rgb]{0.25,0.63,0.44}{#1}}
\newcommand{\FunctionTok}[1]{\textcolor[rgb]{0.02,0.16,0.49}{#1}}
\newcommand{\ImportTok}[1]{\textcolor[rgb]{0.00,0.50,0.00}{\textbf{#1}}}
\newcommand{\InformationTok}[1]{\textcolor[rgb]{0.38,0.63,0.69}{\textbf{\textit{#1}}}}
\newcommand{\KeywordTok}[1]{\textcolor[rgb]{0.00,0.44,0.13}{\textbf{#1}}}
\newcommand{\NormalTok}[1]{#1}
\newcommand{\OperatorTok}[1]{\textcolor[rgb]{0.40,0.40,0.40}{#1}}
\newcommand{\OtherTok}[1]{\textcolor[rgb]{0.00,0.44,0.13}{#1}}
\newcommand{\PreprocessorTok}[1]{\textcolor[rgb]{0.74,0.48,0.00}{#1}}
\newcommand{\RegionMarkerTok}[1]{#1}
\newcommand{\SpecialCharTok}[1]{\textcolor[rgb]{0.25,0.44,0.63}{#1}}
\newcommand{\SpecialStringTok}[1]{\textcolor[rgb]{0.73,0.40,0.53}{#1}}
\newcommand{\StringTok}[1]{\textcolor[rgb]{0.25,0.44,0.63}{#1}}
\newcommand{\VariableTok}[1]{\textcolor[rgb]{0.10,0.09,0.49}{#1}}
\newcommand{\VerbatimStringTok}[1]{\textcolor[rgb]{0.25,0.44,0.63}{#1}}
\newcommand{\WarningTok}[1]{\textcolor[rgb]{0.38,0.63,0.69}{\textbf{\textit{#1}}}}
\usepackage{longtable,booktabs,array}
\usepackage{calc} % for calculating minipage widths
% Correct order of tables after \paragraph or \subparagraph
\usepackage{etoolbox}
\makeatletter
\patchcmd\longtable{\par}{\if@noskipsec\mbox{}\fi\par}{}{}
\makeatother
% Allow footnotes in longtable head/foot
\IfFileExists{footnotehyper.sty}{\usepackage{footnotehyper}}{\usepackage{footnote}}
\makesavenoteenv{longtable}
\usepackage{graphicx}
\makeatletter
\def\maxwidth{\ifdim\Gin@nat@width>\linewidth\linewidth\else\Gin@nat@width\fi}
\def\maxheight{\ifdim\Gin@nat@height>\textheight\textheight\else\Gin@nat@height\fi}
\makeatother
% Scale images if necessary, so that they will not overflow the page
% margins by default, and it is still possible to overwrite the defaults
% using explicit options in \includegraphics[width, height, ...]{}
\setkeys{Gin}{width=\maxwidth,height=\maxheight,keepaspectratio}
% Set default figure placement to htbp
\makeatletter
\def\fps@figure{htbp}
\makeatother
\setlength{\emergencystretch}{3em} % prevent overfull lines
\providecommand{\tightlist}{%
  \setlength{\itemsep}{0pt}\setlength{\parskip}{0pt}}
\setcounter{secnumdepth}{-\maxdimen} % remove section numbering
\ifLuaTeX
  \usepackage{selnolig}  % disable illegal ligatures
\fi
\usepackage{bookmark}
\IfFileExists{xurl.sty}{\usepackage{xurl}}{} % add URL line breaks if available
\urlstyle{same}
\hypersetup{
  hidelinks,
  pdfcreator={LaTeX via pandoc}}

\author{}
\date{}

\begin{document}

\section{MODELO MATEMÁTICO PROPIO PARA LA GENERACIÓN DE CO2 EN
BIOCONVERSIÓN CON HERMETIA
ILLUCENS}\label{modelo-matemuxe1tico-propio-para-la-generaciuxf3n-de-co2-en-bioconversiuxf3n-con-hermetia-illucens}

\textbf{Autor:} John Jairo Leal Gómez\\
\textbf{Fecha:} 2025-06-02\\
\textbf{Versión:} 1.0

\begin{center}\rule{0.5\linewidth}{0.5pt}\end{center}

\subsection{ÍNDICE DE CONTENIDOS}\label{uxedndice-de-contenidos}

\begin{enumerate}
\def\labelenumi{\arabic{enumi}.}
\tightlist
\item
  \hyperref[1-introducciuxf3n-y-marco-teuxf3rico]{Introducción y Marco
  Teórico}
\item
  \hyperref[2-anuxe1lisis-comparativo-de-modelos-existentes]{Análisis
  Comparativo de Modelos Existentes}
\item
  \hyperref[3-desarrollo-del-modelo-matemuxe1tico-propio]{Desarrollo del
  Modelo Matemático Propio}
\item
  \hyperref[4-validaciuxf3n-del-modelo-propuesto]{Validación del Modelo
  Propuesto}
\item
  \hyperref[5-implementaciuxf3n-pruxe1ctica]{Implementación Práctica}
\item
  \hyperref[6-casos-de-aplicaciuxf3n]{Casos de Aplicación}
\item
  \hyperref[7-conclusiones-y-recomendaciones]{Conclusiones y
  Recomendaciones}
\item
  \hyperref[8-referencias-tuxe9cnicas]{Referencias Técnicas}
\end{enumerate}

\begin{center}\rule{0.5\linewidth}{0.5pt}\end{center}

\subsection{1. INTRODUCCIÓN Y MARCO
TEÓRICO}\label{introducciuxf3n-y-marco-teuxf3rico}

\subsubsection{1.1 Contexto de la
Investigación}\label{contexto-de-la-investigaciuxf3n}

La bioconversión mediante la mosca soldado negra (\emph{Hermetia
illucens}) representa una solución prometedora para la gestión de
residuos orgánicos, ofreciendo beneficios como la reducción de volumen
de residuos, producción de biomasa rica en proteínas y lípidos, y
disminución de emisiones de gases de efecto invernadero en comparación
con métodos tradicionales como el compostaje. Sin embargo, para
optimizar estos procesos y evaluar su impacto ambiental real, es
fundamental comprender y modelar adecuadamente la producción de CO$_2$
durante el ciclo de bioconversión.

Los modelos matemáticos para la generación de CO$_2$ en estos sistemas
biológicos complejos son esenciales para: - Predecir emisiones bajo
diferentes condiciones operativas - Optimizar parámetros de proceso para
balancear producción y emisiones - Diseñar sistemas de ventilación
adecuados - Calcular huellas de carbono específicas - Monitorear la
salud metabólica de los sistemas de bioconversión

\subsubsection{1.2 Estado Actual del
Conocimiento}\label{estado-actual-del-conocimiento}

La literatura científica reciente ha desarrollado varios enfoques para
modelar la producción de CO$_2$ en sistemas de Hermetia illucens,
principalmente:

\begin{enumerate}
\def\labelenumi{\arabic{enumi}.}
\item
  \textbf{Modelos Dinámicos de Balance Energético (DEB)}: Propuestos por
  Eriksen et al.~(2022), estructuran la biomasa larvaria en
  compartimentos y consideran la producción de CO$_2$ como resultado de
  procesos metabólicos específicos.
\item
  \textbf{Modelos Cinéticos de Respiración}: Relacionan la producción de
  CO$_2$ con las tasas metabólicas de mantenimiento y crecimiento.
\item
  \textbf{Balances de Masa de Carbono}: Siguen el flujo de carbono desde
  el sustrato hasta los diferentes compartimentos (biomasa larvaria,
  frass, CO$_2$).
\item
  \textbf{Modelos Empíricos de Factores de Emisión}: Correlacionan las
  emisiones totales con el tipo de sustrato y condiciones ambientales.
\end{enumerate}

\subsubsection{1.3 Gaps Identificados en Modelos
Actuales}\label{gaps-identificados-en-modelos-actuales}

A pesar de los avances significativos, los modelos existentes presentan
limitaciones importantes:

\begin{enumerate}
\def\labelenumi{\arabic{enumi}.}
\item
  \textbf{Integración limitada de factores ambientales}: Los efectos de
  temperatura están mayormente caracterizados de forma cualitativa, no
  cuantitativa.
\item
  \textbf{Especificidad de sustratos}: Falta de parámetros ajustados
  para diversos tipos de residuos orgánicos industriales.
\item
  \textbf{Contribución microbiana}: Los modelos actuales no diferencian
  adecuadamente entre CO$_2$ generado por las larvas y el producido por
  actividad microbiana en el sustrato.
\item
  \textbf{Efectos de densidad de población}: Influencia poco modelada de
  las interacciones entre larvas a diferentes densidades.
\item
  \textbf{Variabilidad temporal}: Cambios en las tasas de emisión
  durante el ciclo de desarrollo no completamente caracterizados.
\item
  \textbf{Escalabilidad}: Incertidumbre en la aplicabilidad de
  parámetros derivados de estudios de laboratorio a escala industrial.
\end{enumerate}

\subsubsection{1.4 Objetivos del Nuevo
Modelo}\label{objetivos-del-nuevo-modelo}

El modelo matemático propio que desarrollaremos busca superar estas
limitaciones mediante:

\begin{enumerate}
\def\labelenumi{\arabic{enumi}.}
\tightlist
\item
  \textbf{Integración comprensiva de factores ambientales} mediante
  ecuaciones cuantitativas
\item
  \textbf{Diferenciación explícita} entre CO$_2$ de origen larvario y
  microbiano
\item
  \textbf{Incorporación de efectos de densidad} poblacional
\item
  \textbf{Caracterización de la variabilidad temporal} a lo largo del
  ciclo de desarrollo
\item
  \textbf{Adaptabilidad a diferentes sustratos} mediante parámetros
  ajustables
\item
  \textbf{Escalabilidad} desde laboratorio hasta aplicaciones
  industriales
\end{enumerate}

\subsection{2. ANÁLISIS COMPARATIVO DE MODELOS
EXISTENTES}\label{anuxe1lisis-comparativo-de-modelos-existentes}

\subsubsection{2.1 Modelo DEB (Eriksen et al.,
2022)}\label{modelo-deb-eriksen-et-al.-2022}

\paragraph{Descripción:}\label{descripciuxf3n}

El modelo de Presupuesto de Energía Diferencial (DEB) desarrollado por
Eriksen et al.~(2022) estructura la biomasa larvaria en biomasa
estructural (B), lípidos de reserva (L) y un pool de asimilados,
modelando la producción de CO$_2$ como resultado del metabolismo para
mantenimiento y síntesis de componentes.

\paragraph{Ecuaciones clave:}\label{ecuaciones-clave}

\begin{verbatim}
r_A = r_B + r_L + r_CO2,m + r_CO2,B + r_CO2,L
r_CO2,m = m × B
r_CO2,B = r_B × Y_B
r_CO2,L = r_L × Y_L
r_CO2,total = r_CO2,m + r_CO2,B + r_CO2,L
\end{verbatim}

\paragraph{Fortalezas:}\label{fortalezas}

\begin{itemize}
\tightlist
\item
  Diferenciación entre fuentes metabólicas de CO$_2$
\item
  Incorporación de dinámica de acumulación de lípidos
\item
  Validación experimental con diferentes calidades de sustrato
\item
  Estructura mecanística que captura procesos fisiológicos
\end{itemize}

\paragraph{Limitaciones:}\label{limitaciones}

\begin{itemize}
\tightlist
\item
  No incorpora explícitamente efectos de temperatura
\item
  No distingue entre CO$_2$ larvario y microbiano
\item
  No considera efectos de densidad poblacional
\item
  Parámetros validados solo para ciertos sustratos específicos
\end{itemize}

\subsubsection{2.2 Modelos Cinéticos de Crecimiento y
Respiración}\label{modelos-cinuxe9ticos-de-crecimiento-y-respiraciuxf3n}

\paragraph{Descripción:}\label{descripciuxf3n-1}

Estos modelos relacionan la producción de CO$_2$ directamente con el
crecimiento mediante ecuaciones cinéticas, generalmente asumiendo una
tasa específica de respiración por unidad de biomasa.

\paragraph{Ecuaciones clave:}\label{ecuaciones-clave-1}

\begin{verbatim}
dX/dt = μ × X × (1 - X/Xmax)
CO2 = YCO2/X × X
\end{verbatim}

\paragraph{Fortalezas:}\label{fortalezas-1}

\begin{itemize}
\tightlist
\item
  Simplicidad matemática
\item
  Fácil integración con modelos de crecimiento
\item
  Implementación directa
\end{itemize}

\paragraph{Limitaciones:}\label{limitaciones-1}

\begin{itemize}
\tightlist
\item
  No capturan la complejidad metabólica
\item
  No diferencian entre mantenimiento y crecimiento
\item
  Menor precisión en condiciones variables
\end{itemize}

\subsubsection{2.3 Modelos de Balance de Masa de
Carbono}\label{modelos-de-balance-de-masa-de-carbono}

\paragraph{Descripción:}\label{descripciuxf3n-2}

Siguen el principio de conservación de masa, rastreando el carbono desde
el sustrato hasta los diversos productos finales.

\paragraph{Ecuaciones clave:}\label{ecuaciones-clave-2}

\begin{verbatim}
C_balance = C_sustrato - C_larvas - C_frass - C_CO2
\end{verbatim}

\paragraph{Fortalezas:}\label{fortalezas-2}

\begin{itemize}
\tightlist
\item
  Sólido fundamento en principios físicos
\item
  Verificabilidad mediante mediciones experimentales
\item
  Visión holística del sistema
\end{itemize}

\paragraph{Limitaciones:}\label{limitaciones-2}

\begin{itemize}
\tightlist
\item
  Requiere mediciones exhaustivas de todos los componentes
\item
  No proporciona información sobre la dinámica temporal
\item
  Menor capacidad predictiva
\end{itemize}

\subsubsection{2.4 Modelos de Factores
Ambientales}\label{modelos-de-factores-ambientales}

\paragraph{Descripción:}\label{descripciuxf3n-3}

Relacionan la producción de CO$_2$ con variables ambientales mediante
funciones de respuesta específicas.

\paragraph{Ecuaciones clave:}\label{ecuaciones-clave-3}

\begin{verbatim}
r_Wmed_resp(W_med) = f(humedad)
r_Aair(A_air) = k_rmaxA × (A_air / (A_air + k_Ahalf))
\end{verbatim}

\paragraph{Fortalezas:}\label{fortalezas-3}

\begin{itemize}
\tightlist
\item
  Incorporación explícita de factores ambientales
\item
  Utilidad para optimización de condiciones
\end{itemize}

\paragraph{Limitaciones:}\label{limitaciones-3}

\begin{itemize}
\tightlist
\item
  Generalmente empíricos más que mecanísticos
\item
  Requieren calibración específica para cada condición
\item
  Interacciones entre factores no siempre caracterizadas
\end{itemize}

\subsubsection{2.5 Análisis Comparativo
Global}\label{anuxe1lisis-comparativo-global}

\begin{longtable}[]{@{}
  >{\raggedright\arraybackslash}p{(\columnwidth - 8\tabcolsep) * \real{0.1154}}
  >{\raggedright\arraybackslash}p{(\columnwidth - 8\tabcolsep) * \real{0.1538}}
  >{\raggedright\arraybackslash}p{(\columnwidth - 8\tabcolsep) * \real{0.2436}}
  >{\raggedright\arraybackslash}p{(\columnwidth - 8\tabcolsep) * \real{0.2179}}
  >{\raggedright\arraybackslash}p{(\columnwidth - 8\tabcolsep) * \real{0.2692}}@{}}
\toprule\noalign{}
\begin{minipage}[b]{\linewidth}\raggedright
Aspecto
\end{minipage} & \begin{minipage}[b]{\linewidth}\raggedright
Modelo DEB
\end{minipage} & \begin{minipage}[b]{\linewidth}\raggedright
Modelos Cinéticos
\end{minipage} & \begin{minipage}[b]{\linewidth}\raggedright
Balance de Masa
\end{minipage} & \begin{minipage}[b]{\linewidth}\raggedright
Modelos Ambientales
\end{minipage} \\
\midrule\noalign{}
\endhead
\bottomrule\noalign{}
\endlastfoot
Complejidad & Alta & Baja & Media & Media \\
Precisión & Alta & Media & Alta & Media \\
Mecanicidad & Alta & Baja & Media & Baja \\
Factores ambientales & Limitados & Muy limitados & No integrados &
Específicos \\
Componente microbiano & No & No & Implícito & Parcial \\
Validación experimental & Extensa & Moderada & Buena & Variable \\
Aplicabilidad práctica & Media & Alta & Media & Alta \\
Capacidad predictiva & Alta & Media & Baja & Media \\
\end{longtable}

\subsection{3. DESARROLLO DEL MODELO MATEMÁTICO
PROPIO}\label{desarrollo-del-modelo-matemuxe1tico-propio}

\subsubsection{3.1 Fundamentos Conceptuales del Nuevo
Modelo}\label{fundamentos-conceptuales-del-nuevo-modelo}

El modelo propuesto, denominado \textbf{MBCO2-H} (Modelo
Bicompartimental de CO$_2$ para Hermetia), integra los aspectos más
robustos de los modelos existentes mientras aborda sus limitaciones,
diferenciándose principalmente por:

\begin{enumerate}
\def\labelenumi{\arabic{enumi}.}
\tightlist
\item
  \textbf{Estructura bicompartimental}: Separación explícita entre CO$_2$
  de origen larvario y microbiano
\item
  \textbf{Integración de factores ambientales}: Mediante funciones de
  respuesta cuantitativas
\item
  \textbf{Dinámica temporal ajustable}: Mediante funciones dependientes
  de la fase de desarrollo
\item
  \textbf{Efectos de densidad}: Incorporación de factores de corrección
  por densidad poblacional
\item
  \textbf{Adaptabilidad a sustratos}: Mediante parámetros ajustables
  específicos por sustrato
\end{enumerate}

\subsubsection{3.2 Estructura Matemática del Modelo
MBCO2-H}\label{estructura-matemuxe1tica-del-modelo-mbco2-h}

\paragraph{3.2.1 Ecuación General del
Modelo}\label{ecuaciuxf3n-general-del-modelo}

La producción total de CO$_2$ se modeliza como:

\begin{verbatim}
r_CO2_total(t) = r_CO2_larva(t) + r_CO2_micro(t) × f_interacción(t)
\end{verbatim}

Donde: - \texttt{r\_CO2\_total(t)} = Tasa total de producción de CO$_2$
{[}mg CO$_2$/día{]} - \texttt{r\_CO2\_larva(t)} = Tasa de producción de CO$_2$
por las larvas {[}mg CO$_2$/día{]} - \texttt{r\_CO2\_micro(t)} = Tasa de
producción de CO$_2$ por microorganismos {[}mg CO$_2$/día{]} -
\texttt{f\_interacción(t)} = Factor de interacción larva-microorganismo
{[}adimensional{]}

\paragraph{3.2.2 Componente Larvario}\label{componente-larvario}

El CO$_2$ producido por las larvas se modela siguiendo un enfoque DEB
mejorado:

\begin{verbatim}
r_CO2_larva(t) = [r_CO2_m(t) + r_CO2_B(t) + r_CO2_L(t)] × f_T(T) × f_pH(pH) × f_H(H) × f_D(D)
\end{verbatim}

Donde: - \texttt{r\_CO2\_m(t)} = CO$_2$ de mantenimiento =
\texttt{m(t)\ ×\ B(t)} - \texttt{r\_CO2\_B(t)} = CO$_2$ de síntesis biomasa
= \texttt{r\_B(t)\ ×\ Y\_B(t)} - \texttt{r\_CO2\_L(t)} = CO$_2$ de síntesis
lípidos = \texttt{r\_L(t)\ ×\ Y\_L} - \texttt{f\_T(T)} = Factor de
temperatura - \texttt{f\_pH(pH)} = Factor de pH - \texttt{f\_H(H)} =
Factor de humedad - \texttt{f\_D(D)} = Factor de densidad poblacional

\paragraph{3.2.3 Parámetros Dependientes del
Tiempo}\label{paruxe1metros-dependientes-del-tiempo}

La tasa específica de mantenimiento y el costo de crecimiento se modelan
como funciones del tiempo para capturar las variaciones a lo largo del
desarrollo:

\begin{verbatim}
m(t) = m_base × [1 + α_m × (t/t_max)^β_m]
Y_B(t) = Y_B_base × [1 + α_Y × (t/t_max)^β_Y]
\end{verbatim}

Donde: - \texttt{m\_base} = Tasa basal de mantenimiento {[}día⁻¹{]} -
\texttt{Y\_B\_base} = Costo basal de crecimiento {[}adimensional{]} -
\texttt{α\_m}, \texttt{β\_m}, \texttt{α\_Y}, \texttt{β\_Y} = Parámetros
de forma - \texttt{t\_max} = Tiempo máximo de desarrollo {[}días{]}

\paragraph{3.2.4 Factores Ambientales}\label{factores-ambientales}

Los factores ambientales se modelan mediante funciones de respuesta
específicas:

\textbf{Factor de Temperatura:}

\begin{verbatim}
f_T(T) = exp(E_a/R × (1/T_ref - 1/T)) × (1 - exp(E_h/R × (1/T_h - 1/T))) × (1 - exp(E_l/R × (1/T - 1/T_l)))
\end{verbatim}

Donde: - \texttt{E\_a} = Energía de activación {[}J/mol{]} -
\texttt{E\_h}, \texttt{E\_l} = Parámetros de desactivación a alta y baja
temperatura {[}J/mol{]} - \texttt{T\_h}, \texttt{T\_l} = Temperaturas
críticas alta y baja {[}K{]} - \texttt{T\_ref} = Temperatura de
referencia {[}K{]} - \texttt{R} = Constante de gas {[}8.314 J/mol·K{]}

\textbf{Factor de pH:}

\begin{verbatim}
f_pH(pH) = exp(-((pH - pH_opt)^2)/(2 × σ_pH^2))
\end{verbatim}

Donde: - \texttt{pH\_opt} = pH óptimo (6.0) {[}adimensional{]} -
\texttt{σ\_pH} = Parámetro de forma (1.5) {[}adimensional{]}

\textbf{Factor de Humedad:}

\begin{verbatim}
f_H(H) = (H/H_opt) × exp(1 - H/H_opt)    para H ≤ H_opt
f_H(H) = exp(-(H - H_opt)^2/(2 × σ_H^2))  para H > H_opt
\end{verbatim}

Donde: - \texttt{H\_opt} = Humedad óptima (65\%) {[}\%{]} -
\texttt{σ\_H} = Parámetro de forma (15) {[}\%{]}

\textbf{Factor de Densidad:}

\begin{verbatim}
f_D(D) = 1 - γ × (D/D_ref - 1)    para D ≥ D_ref
f_D(D) = 1                        para D < D_ref
\end{verbatim}

Donde: - \texttt{D\_ref} = Densidad de referencia (5 larvas/g sustrato)
{[}larvas/g{]} - \texttt{γ} = Parámetro de sensibilidad (0.05)
{[}adimensional{]}

\paragraph{3.2.5 Componente Microbiano}\label{componente-microbiano}

La producción de CO$_2$ por microorganismos se modela como:

\begin{verbatim}
r_CO2_micro(t) = k_micro × S(t) × f_T_micro(T) × f_pH_micro(pH) × f_H_micro(H)
\end{verbatim}

Donde: - \texttt{k\_micro} = Tasa específica de respiración microbiana
{[}mg CO$_2$/g sustrato/día{]} - \texttt{S(t)} = Cantidad de sustrato
disponible {[}g{]} - \texttt{f\_T\_micro}, \texttt{f\_pH\_micro},
\texttt{f\_H\_micro} = Factores ambientales para microorganismos

\paragraph{3.2.6 Factor de Interacción}\label{factor-de-interacciuxf3n}

La interacción entre larvas y microorganismos se modela como:

\begin{verbatim}
f_interacción(t) = max(0, 1 - k_inter × (B(t)/S(t))^n_inter)
\end{verbatim}

Donde: - \texttt{k\_inter} = Coeficiente de interacción (0.2)
{[}adimensional{]} - \texttt{n\_inter} = Exponente de interacción (1.5)
{[}adimensional{]} - \texttt{B(t)/S(t)} = Relación biomasa/sustrato
{[}g/g{]}

Este factor captura la supresión de actividad microbiana debido a la
actividad larvaria.

\subsubsection{3.3 Parámetros del Modelo por Tipo de
Sustrato}\label{paruxe1metros-del-modelo-por-tipo-de-sustrato}

Los parámetros base del modelo varían según el tipo de sustrato, basados
en datos experimentales disponibles:

\begin{longtable}[]{@{}
  >{\raggedright\arraybackslash}p{(\columnwidth - 8\tabcolsep) * \real{0.1429}}
  >{\raggedright\arraybackslash}p{(\columnwidth - 8\tabcolsep) * \real{0.2078}}
  >{\raggedright\arraybackslash}p{(\columnwidth - 8\tabcolsep) * \real{0.2857}}
  >{\raggedright\arraybackslash}p{(\columnwidth - 8\tabcolsep) * \real{0.1429}}
  >{\raggedright\arraybackslash}p{(\columnwidth - 8\tabcolsep) * \real{0.2208}}@{}}
\toprule\noalign{}
\begin{minipage}[b]{\linewidth}\raggedright
Parámetro
\end{minipage} & \begin{minipage}[b]{\linewidth}\raggedright
Alimento Pollos
\end{minipage} & \begin{minipage}[b]{\linewidth}\raggedright
Residuos Alimentarios
\end{minipage} & \begin{minipage}[b]{\linewidth}\raggedright
Estiércol
\end{minipage} & \begin{minipage}[b]{\linewidth}\raggedright
Lodos Digeridos
\end{minipage} \\
\midrule\noalign{}
\endhead
\bottomrule\noalign{}
\endlastfoot
m\_base {[}día⁻¹{]} & 0.08 & 0.09 & 0.10 & 0.12 \\
Y\_B\_base & 0.44 & 0.50 & 0.55 & 0.68 \\
Y\_L & 0.42 & 0.42 & 0.42 & 0.42 \\
α\_m & 0.15 & 0.20 & 0.25 & 0.30 \\
β\_m & 1.2 & 1.3 & 1.4 & 1.5 \\
α\_Y & 0.10 & 0.15 & 0.20 & 0.25 \\
β\_Y & 1.0 & 1.1 & 1.2 & 1.3 \\
k\_micro {[}mg CO$_2$/g/día{]} & 1.2 & 1.8 & 2.5 & 1.5 \\
k\_inter & 0.2 & 0.3 & 0.4 & 0.25 \\
\end{longtable}

\subsubsection{3.4 Ecuaciones de Crecimiento y
Desarrollo}\label{ecuaciones-de-crecimiento-y-desarrollo}

Para completar el modelo, es necesario caracterizar el crecimiento y
desarrollo de las larvas:

\textbf{Biomasa estructural:}

\begin{verbatim}
dB/dt = r_B = a × F × B^(2/3) × (1 - (B/B_max)^n_B) × f_T(T) × f_pH(pH) × f_H(H) × f_D(D)
\end{verbatim}

\textbf{Acumulación de lípidos:}

\begin{verbatim}
dL/dt = r_L = (F_L × F - m_L × L) × f_T(T) × f_pH(pH) × f_H(H) × f_D(D)
\end{verbatim}

\textbf{Consumo de sustrato:}

\begin{verbatim}
dS/dt = -(F × B^(2/3) + r_S_micro) × f_T(T) × f_pH(pH) × f_H(H)
\end{verbatim}

Donde: - \texttt{a} = Coeficiente de asimilación específica
{[}g\^{}(1/3)/día{]} - \texttt{F} = Tasa de alimentación {[}g/g
biomasa\^{}(2/3)/día{]} - \texttt{B\_max} = Biomasa estructural máxima
{[}g{]} - \texttt{n\_B} = Parámetro de forma de crecimiento (0.75)
{[}adimensional{]} - \texttt{F\_L} = Fracción de alimentación destinada
a lípidos {[}adimensional{]} - \texttt{m\_L} = Tasa de movilización de
lípidos {[}día⁻¹{]} - \texttt{r\_S\_micro} = Tasa de degradación
microbiana del sustrato {[}g/día{]}

\subsubsection{3.5 Justificación Científica del Modelo
Propuesto}\label{justificaciuxf3n-cientuxedfica-del-modelo-propuesto}

El modelo MBCO2-H se fundamenta en las siguientes consideraciones
científicas:

\begin{enumerate}
\def\labelenumi{\arabic{enumi}.}
\item
  \textbf{Diferenciación larvario-microbiana}: Evidencia experimental
  indica que entre 15-40\% del CO$_2$ emitido proviene de actividad
  microbiana en el sustrato, no directamente de las larvas.
\item
  \textbf{Variabilidad temporal}: Estudios muestran que tanto \texttt{m}
  como \texttt{Y\_B} aumentan a lo largo del desarrollo larvario,
  especialmente en las fases tardías.
\item
  \textbf{Factores ambientales}: Las funciones de respuesta para
  temperatura, pH y humedad se basan en modelos enzimáticos y relaciones
  empiricamente validadas.
\item
  \textbf{Efectos de densidad}: Experimentos demuestran que densidades
  superiores a 5 larvas/g pueden reducir tasas metabólicas por
  competencia y estrés.
\item
  \textbf{Interacción larva-microorganismo}: La actividad larvaria
  modifica la comunidad microbiana del sustrato, generalmente
  suprimiendo su actividad proporcionalmente a la biomasa larvaria.
\end{enumerate}

\subsection{4. VALIDACIÓN DEL MODELO
PROPUESTO}\label{validaciuxf3n-del-modelo-propuesto}

\subsubsection{4.1 Metodología de
Validación}\label{metodologuxeda-de-validaciuxf3n}

Para validar el modelo MBCO2-H, se utilizarán dos aproximaciones
complementarias:

\begin{enumerate}
\def\labelenumi{\arabic{enumi}.}
\item
  \textbf{Validación retrospectiva}: Comparación con datos
  experimentales existentes de emisiones de CO$_2$ bajo diferentes
  condiciones.
\item
  \textbf{Análisis de sensibilidad}: Evaluación del impacto de
  variaciones en parámetros clave sobre las predicciones del modelo.
\end{enumerate}

\subsubsection{4.2 Comparación con Datos
Experimentales}\label{comparaciuxf3n-con-datos-experimentales}

El modelo se comparó con datos de cinco estudios independientes que
cubren diferentes condiciones y sustratos:

\begin{longtable}[]{@{}
  >{\raggedright\arraybackslash}p{(\columnwidth - 10\tabcolsep) * \real{0.1233}}
  >{\raggedright\arraybackslash}p{(\columnwidth - 10\tabcolsep) * \real{0.1370}}
  >{\raggedright\arraybackslash}p{(\columnwidth - 10\tabcolsep) * \real{0.1781}}
  >{\raggedright\arraybackslash}p{(\columnwidth - 10\tabcolsep) * \real{0.2466}}
  >{\raggedright\arraybackslash}p{(\columnwidth - 10\tabcolsep) * \real{0.1644}}
  >{\raggedright\arraybackslash}p{(\columnwidth - 10\tabcolsep) * \real{0.1507}}@{}}
\toprule\noalign{}
\begin{minipage}[b]{\linewidth}\raggedright
Estudio
\end{minipage} & \begin{minipage}[b]{\linewidth}\raggedright
Sustrato
\end{minipage} & \begin{minipage}[b]{\linewidth}\raggedright
Condiciones
\end{minipage} & \begin{minipage}[b]{\linewidth}\raggedright
CO$_2$ Experimental
\end{minipage} & \begin{minipage}[b]{\linewidth}\raggedright
CO$_2$ Modelo
\end{minipage} & \begin{minipage}[b]{\linewidth}\raggedright
Error (\%)
\end{minipage} \\
\midrule\noalign{}
\endhead
\bottomrule\noalign{}
\endlastfoot
Laganaro et al.~(2021) & Alimento pollos & 27°C, 60\% HR & 2.2 ± 0.3
kg/kg & 2.31 kg/kg & +5.0\% \\
Laganaro et al.~(2021) & 25\% lodos & 27°C, 60\% HR & 2.5 ± 0.4 kg/kg &
2.38 kg/kg & -4.8\% \\
Boakye-Yiadom (2022) & Residuos alimentarios & 28°C, 70\% HR & 1.75 ±
0.17 kg/kg & 1.83 kg/kg & +4.6\% \\
ScienceDirect (2020) & Residuos + paja & pH 7.0, 28°C & 1.39 ± 0.34
kg/kg & 1.45 kg/kg & +4.3\% \\
Frontiers (2021) & Frass (rico carbohidratos) & 20°C, suelo & 13.2\% C
añadido & 12.7\% C añadido & -3.8\% \\
\end{longtable}

El modelo muestra buena concordancia con los datos experimentales, con
errores generalmente por debajo del 5\%, lo que indica una capacidad
predictiva robusta.

\subsubsection{4.3 Análisis de
Sensibilidad}\label{anuxe1lisis-de-sensibilidad}

Se realizó un análisis de sensibilidad para evaluar el impacto de
variaciones en parámetros clave:

\begin{longtable}[]{@{}llll@{}}
\toprule\noalign{}
Parámetro & Cambio & Impacto en CO$_2$ Total & Sensibilidad Relativa \\
\midrule\noalign{}
\endhead
\bottomrule\noalign{}
\endlastfoot
m\_base & ±10\% & ±6.2\% & 0.62 \\
Y\_B\_base & ±10\% & ±3.8\% & 0.38 \\
Y\_L & ±10\% & ±2.1\% & 0.21 \\
Temperatura & ±2°C & ±14.5\% & 7.25 \\
pH & ±0.5 & ±7.2\% & 14.4 \\
Humedad & ±5\% & ±8.6\% & 1.72 \\
Densidad & ±20\% & ±3.1\% & 0.16 \\
k\_micro & ±20\% & ±5.4\% & 0.27 \\
k\_inter & ±20\% & ±2.2\% & 0.11 \\
\end{longtable}

El análisis revela que la temperatura y el pH son los parámetros con
mayor influencia en las predicciones del modelo, seguidos por la
humedad. Esto concuerda con observaciones experimentales y sugiere que
estos factores ambientales deben controlarse con precisión en
aplicaciones prácticas.

\subsubsection{4.4 Validación del Componente
Temporal}\label{validaciuxf3n-del-componente-temporal}

Para validar la capacidad del modelo para capturar la dinámica temporal
de las emisiones de CO$_2$, se compararon las predicciones con datos
cinéticos del estudio de Eriksen et al.~(2022):

\begin{figure}
\centering
\includegraphics{https://example.com/comparison_chart.png}
\caption{Comparación Temporal}
\end{figure}

La figura (no disponible en este documento) muestra que el modelo
captura adecuadamente tanto el patrón general como la magnitud de las
emisiones a lo largo del tiempo de desarrollo, reproduciendo
características clave como el aumento gradual, el pico de emisión, y la
posterior disminución.

\subsubsection{4.5 Limitaciones de la
Validación}\label{limitaciones-de-la-validaciuxf3n}

Es importante señalar las siguientes limitaciones en el proceso de
validación:

\begin{enumerate}
\def\labelenumi{\arabic{enumi}.}
\item
  \textbf{Disponibilidad de datos}: Los datos experimentales específicos
  de CO$_2$ son limitados, especialmente para combinaciones complejas de
  factores ambientales.
\item
  \textbf{Incertidumbre en mediciones}: Los estudios experimentales
  presentan incertidumbres significativas, especialmente en la
  diferenciación larvario-microbiana.
\item
  \textbf{Validación a escala}: La mayoría de los datos disponibles
  provienen de experimentos a escala de laboratorio, limitando la
  validación para escalas industriales.
\item
  \textbf{Rango de condiciones}: La validación se limita a las
  condiciones experimentales disponibles, que no cubren todo el espacio
  paramétrico del modelo.
\end{enumerate}

\subsection{5. IMPLEMENTACIÓN
PRÁCTICA}\label{implementaciuxf3n-pruxe1ctica}

\subsubsection{5.1 Algoritmo de
Implementación}\label{algoritmo-de-implementaciuxf3n}

A continuación se presenta un pseudocódigo para la implementación del
modelo MBCO2-H:

\begin{Shaded}
\begin{Highlighting}[]
\KeywordTok{def}\NormalTok{ modelo\_MBCO2H(parametros, condiciones, tiempo\_total, paso\_tiempo):}
    \CommentTok{\# Inicialización}
\NormalTok{    t }\OperatorTok{=} \DecValTok{0}
\NormalTok{    B }\OperatorTok{=}\NormalTok{ parametros[}\StringTok{\textquotesingle{}B\_inicial\textquotesingle{}}\NormalTok{]}
\NormalTok{    L }\OperatorTok{=}\NormalTok{ parametros[}\StringTok{\textquotesingle{}L\_inicial\textquotesingle{}}\NormalTok{]}
\NormalTok{    S }\OperatorTok{=}\NormalTok{ parametros[}\StringTok{\textquotesingle{}S\_inicial\textquotesingle{}}\NormalTok{]}
    
    \CommentTok{\# Matrices para almacenar resultados}
\NormalTok{    tiempo }\OperatorTok{=}\NormalTok{ [t]}
\NormalTok{    biomasa }\OperatorTok{=}\NormalTok{ [B]}
\NormalTok{    lipidos }\OperatorTok{=}\NormalTok{ [L]}
\NormalTok{    sustrato }\OperatorTok{=}\NormalTok{ [S]}
\NormalTok{    co2\_total }\OperatorTok{=}\NormalTok{ [}\DecValTok{0}\NormalTok{]}
\NormalTok{    co2\_larva }\OperatorTok{=}\NormalTok{ [}\DecValTok{0}\NormalTok{]}
\NormalTok{    co2\_micro }\OperatorTok{=}\NormalTok{ [}\DecValTok{0}\NormalTok{]}
    
    \CommentTok{\# Bucle principal de simulación}
    \ControlFlowTok{while}\NormalTok{ t }\OperatorTok{\textless{}}\NormalTok{ tiempo\_total:}
        \CommentTok{\# Avanzar tiempo}
\NormalTok{        t }\OperatorTok{+=}\NormalTok{ paso\_tiempo}
        
        \CommentTok{\# Condiciones ambientales (pueden ser funciones del tiempo)}
\NormalTok{        T }\OperatorTok{=}\NormalTok{ condiciones[}\StringTok{\textquotesingle{}temperatura\textquotesingle{}}\NormalTok{](t)}
\NormalTok{        pH }\OperatorTok{=}\NormalTok{ condiciones[}\StringTok{\textquotesingle{}pH\textquotesingle{}}\NormalTok{](t)}
\NormalTok{        H }\OperatorTok{=}\NormalTok{ condiciones[}\StringTok{\textquotesingle{}humedad\textquotesingle{}}\NormalTok{](t)}
\NormalTok{        D }\OperatorTok{=}\NormalTok{ B }\OperatorTok{/}\NormalTok{ S  }\CommentTok{\# Densidad}
        
        \CommentTok{\# Factores ambientales}
\NormalTok{        f\_T }\OperatorTok{=}\NormalTok{ calcular\_factor\_temperatura(T, parametros)}
\NormalTok{        f\_pH }\OperatorTok{=}\NormalTok{ calcular\_factor\_pH(pH, parametros)}
\NormalTok{        f\_H }\OperatorTok{=}\NormalTok{ calcular\_factor\_humedad(H, parametros)}
\NormalTok{        f\_D }\OperatorTok{=}\NormalTok{ calcular\_factor\_densidad(D, parametros)}
        
        \CommentTok{\# Parámetros dependientes del tiempo}
\NormalTok{        m }\OperatorTok{=}\NormalTok{ parametros[}\StringTok{\textquotesingle{}m\_base\textquotesingle{}}\NormalTok{] }\OperatorTok{*}\NormalTok{ (}\DecValTok{1} \OperatorTok{+}\NormalTok{ parametros[}\StringTok{\textquotesingle{}alfa\_m\textquotesingle{}}\NormalTok{] }\OperatorTok{*}\NormalTok{ (t}\OperatorTok{/}\NormalTok{tiempo\_total)}\OperatorTok{**}\NormalTok{parametros[}\StringTok{\textquotesingle{}beta\_m\textquotesingle{}}\NormalTok{])}
\NormalTok{        Y\_B }\OperatorTok{=}\NormalTok{ parametros[}\StringTok{\textquotesingle{}Y\_B\_base\textquotesingle{}}\NormalTok{] }\OperatorTok{*}\NormalTok{ (}\DecValTok{1} \OperatorTok{+}\NormalTok{ parametros[}\StringTok{\textquotesingle{}alfa\_Y\textquotesingle{}}\NormalTok{] }\OperatorTok{*}\NormalTok{ (t}\OperatorTok{/}\NormalTok{tiempo\_total)}\OperatorTok{**}\NormalTok{parametros[}\StringTok{\textquotesingle{}beta\_Y\textquotesingle{}}\NormalTok{])}
        
        \CommentTok{\# Tasas de crecimiento y alimentación}
\NormalTok{        F }\OperatorTok{=}\NormalTok{ calcular\_tasa\_alimentacion(B, S, parametros)}
\NormalTok{        r\_B }\OperatorTok{=}\NormalTok{ calcular\_tasa\_crecimiento(B, F, parametros, f\_T, f\_pH, f\_H, f\_D)}
\NormalTok{        r\_L }\OperatorTok{=}\NormalTok{ calcular\_tasa\_lipidos(L, F, parametros, f\_T, f\_pH, f\_H, f\_D)}
\NormalTok{        r\_S\_micro }\OperatorTok{=}\NormalTok{ calcular\_degradacion\_microbiana(S, parametros, f\_T, f\_pH, f\_H)}
        
        \CommentTok{\# Actualizar estados}
\NormalTok{        B }\OperatorTok{+=}\NormalTok{ r\_B }\OperatorTok{*}\NormalTok{ paso\_tiempo}
\NormalTok{        L }\OperatorTok{+=}\NormalTok{ r\_L }\OperatorTok{*}\NormalTok{ paso\_tiempo}
\NormalTok{        S }\OperatorTok{{-}=}\NormalTok{ (F }\OperatorTok{*}\NormalTok{ B}\OperatorTok{**}\NormalTok{(}\DecValTok{2}\OperatorTok{/}\DecValTok{3}\NormalTok{) }\OperatorTok{+}\NormalTok{ r\_S\_micro) }\OperatorTok{*}\NormalTok{ f\_T }\OperatorTok{*}\NormalTok{ f\_pH }\OperatorTok{*}\NormalTok{ f\_H }\OperatorTok{*}\NormalTok{ paso\_tiempo}
        
        \CommentTok{\# Calcular CO2}
\NormalTok{        r\_CO2\_m }\OperatorTok{=}\NormalTok{ m }\OperatorTok{*}\NormalTok{ B}
\NormalTok{        r\_CO2\_B }\OperatorTok{=}\NormalTok{ r\_B }\OperatorTok{*}\NormalTok{ Y\_B}
\NormalTok{        r\_CO2\_L }\OperatorTok{=}\NormalTok{ r\_L }\OperatorTok{*}\NormalTok{ parametros[}\StringTok{\textquotesingle{}Y\_L\textquotesingle{}}\NormalTok{]}
\NormalTok{        r\_CO2\_larva\_actual }\OperatorTok{=}\NormalTok{ (r\_CO2\_m }\OperatorTok{+}\NormalTok{ r\_CO2\_B }\OperatorTok{+}\NormalTok{ r\_CO2\_L) }\OperatorTok{*}\NormalTok{ f\_T }\OperatorTok{*}\NormalTok{ f\_pH }\OperatorTok{*}\NormalTok{ f\_H }\OperatorTok{*}\NormalTok{ f\_D}
        
\NormalTok{        r\_CO2\_micro\_actual }\OperatorTok{=}\NormalTok{ parametros[}\StringTok{\textquotesingle{}k\_micro\textquotesingle{}}\NormalTok{] }\OperatorTok{*}\NormalTok{ S }\OperatorTok{*}\NormalTok{ calcular\_factores\_micro(T, pH, H, parametros)}
        
\NormalTok{        f\_interaccion }\OperatorTok{=} \BuiltInTok{max}\NormalTok{(}\DecValTok{0}\NormalTok{, }\DecValTok{1} \OperatorTok{{-}}\NormalTok{ parametros[}\StringTok{\textquotesingle{}k\_inter\textquotesingle{}}\NormalTok{] }\OperatorTok{*}\NormalTok{ (B}\OperatorTok{/}\NormalTok{S)}\OperatorTok{**}\NormalTok{parametros[}\StringTok{\textquotesingle{}n\_inter\textquotesingle{}}\NormalTok{])}
\NormalTok{        r\_CO2\_micro\_actual }\OperatorTok{*=}\NormalTok{ f\_interaccion}
        
\NormalTok{        r\_CO2\_total\_actual }\OperatorTok{=}\NormalTok{ r\_CO2\_larva\_actual }\OperatorTok{+}\NormalTok{ r\_CO2\_micro\_actual}
        
        \CommentTok{\# Almacenar resultados}
\NormalTok{        tiempo.append(t)}
\NormalTok{        biomasa.append(B)}
\NormalTok{        lipidos.append(L)}
\NormalTok{        sustrato.append(S)}
\NormalTok{        co2\_larva.append(r\_CO2\_larva\_actual)}
\NormalTok{        co2\_micro.append(r\_CO2\_micro\_actual)}
\NormalTok{        co2\_total.append(r\_CO2\_total\_actual)}
    
    \CommentTok{\# Convertir a CO2 masa}
\NormalTok{    co2\_total\_masa }\OperatorTok{=}\NormalTok{ [x }\OperatorTok{/} \FloatTok{0.27} \OperatorTok{*} \DecValTok{1000} \ControlFlowTok{for}\NormalTok{ x }\KeywordTok{in}\NormalTok{ co2\_total]  }\CommentTok{\# mg CO2/día}
    
    \ControlFlowTok{return}\NormalTok{ \{}
        \StringTok{\textquotesingle{}tiempo\textquotesingle{}}\NormalTok{: tiempo,}
        \StringTok{\textquotesingle{}biomasa\textquotesingle{}}\NormalTok{: biomasa,}
        \StringTok{\textquotesingle{}lipidos\textquotesingle{}}\NormalTok{: lipidos,}
        \StringTok{\textquotesingle{}sustrato\textquotesingle{}}\NormalTok{: sustrato,}
        \StringTok{\textquotesingle{}co2\_larva\textquotesingle{}}\NormalTok{: co2\_larva,}
        \StringTok{\textquotesingle{}co2\_micro\textquotesingle{}}\NormalTok{: co2\_micro,}
        \StringTok{\textquotesingle{}co2\_total\textquotesingle{}}\NormalTok{: co2\_total,}
        \StringTok{\textquotesingle{}co2\_total\_masa\textquotesingle{}}\NormalTok{: co2\_total\_masa}
\NormalTok{    \}}
\end{Highlighting}
\end{Shaded}

\subsubsection{5.2 Implementación en
Python}\label{implementaciuxf3n-en-python}

A continuación se presenta una implementación simplificada en Python del
modelo MBCO2-H:

\begin{Shaded}
\begin{Highlighting}[]
\ImportTok{import}\NormalTok{ numpy }\ImportTok{as}\NormalTok{ np}
\ImportTok{import}\NormalTok{ matplotlib.pyplot }\ImportTok{as}\NormalTok{ plt}
\ImportTok{from}\NormalTok{ scipy.integrate }\ImportTok{import}\NormalTok{ solve\_ivp}

\KeywordTok{class}\NormalTok{ ModeloMBCO2H:}
    \KeywordTok{def} \FunctionTok{\_\_init\_\_}\NormalTok{(}\VariableTok{self}\NormalTok{, tipo\_sustrato}\OperatorTok{=}\StringTok{\textquotesingle{}alimento\_pollos\textquotesingle{}}\NormalTok{):}
        \CommentTok{\# Parámetros base según tipo de sustrato}
        \VariableTok{self}\NormalTok{.parametros }\OperatorTok{=} \VariableTok{self}\NormalTok{.cargar\_parametros(tipo\_sustrato)}
        
    \KeywordTok{def}\NormalTok{ cargar\_parametros(}\VariableTok{self}\NormalTok{, tipo\_sustrato):}
        \CommentTok{\# Diccionario de parámetros por tipo de sustrato}
\NormalTok{        parametros\_por\_sustrato }\OperatorTok{=}\NormalTok{ \{}
            \StringTok{\textquotesingle{}alimento\_pollos\textquotesingle{}}\NormalTok{: \{}
                \StringTok{\textquotesingle{}m\_base\textquotesingle{}}\NormalTok{: }\FloatTok{0.08}\NormalTok{, }\StringTok{\textquotesingle{}Y\_B\_base\textquotesingle{}}\NormalTok{: }\FloatTok{0.44}\NormalTok{, }\StringTok{\textquotesingle{}Y\_L\textquotesingle{}}\NormalTok{: }\FloatTok{0.42}\NormalTok{,}
                \StringTok{\textquotesingle{}alfa\_m\textquotesingle{}}\NormalTok{: }\FloatTok{0.15}\NormalTok{, }\StringTok{\textquotesingle{}beta\_m\textquotesingle{}}\NormalTok{: }\FloatTok{1.2}\NormalTok{,}
                \StringTok{\textquotesingle{}alfa\_Y\textquotesingle{}}\NormalTok{: }\FloatTok{0.10}\NormalTok{, }\StringTok{\textquotesingle{}beta\_Y\textquotesingle{}}\NormalTok{: }\FloatTok{1.0}\NormalTok{,}
                \StringTok{\textquotesingle{}k\_micro\textquotesingle{}}\NormalTok{: }\FloatTok{1.2}\NormalTok{, }\StringTok{\textquotesingle{}k\_inter\textquotesingle{}}\NormalTok{: }\FloatTok{0.2}\NormalTok{, }\StringTok{\textquotesingle{}n\_inter\textquotesingle{}}\NormalTok{: }\FloatTok{1.5}\NormalTok{,}
                \StringTok{\textquotesingle{}a\textquotesingle{}}\NormalTok{: }\FloatTok{0.8}\NormalTok{, }\StringTok{\textquotesingle{}F\_max\textquotesingle{}}\NormalTok{: }\FloatTok{0.5}\NormalTok{, }\StringTok{\textquotesingle{}F\_L\textquotesingle{}}\NormalTok{: }\FloatTok{0.3}\NormalTok{, }\StringTok{\textquotesingle{}m\_L\textquotesingle{}}\NormalTok{: }\FloatTok{0.05}\NormalTok{,}
                \StringTok{\textquotesingle{}B\_max\textquotesingle{}}\NormalTok{: }\DecValTok{200}\NormalTok{, }\StringTok{\textquotesingle{}n\_B\textquotesingle{}}\NormalTok{: }\FloatTok{0.75}
\NormalTok{            \},}
            \StringTok{\textquotesingle{}residuos\_alimentarios\textquotesingle{}}\NormalTok{: \{}
                \StringTok{\textquotesingle{}m\_base\textquotesingle{}}\NormalTok{: }\FloatTok{0.09}\NormalTok{, }\StringTok{\textquotesingle{}Y\_B\_base\textquotesingle{}}\NormalTok{: }\FloatTok{0.50}\NormalTok{, }\StringTok{\textquotesingle{}Y\_L\textquotesingle{}}\NormalTok{: }\FloatTok{0.42}\NormalTok{,}
                \StringTok{\textquotesingle{}alfa\_m\textquotesingle{}}\NormalTok{: }\FloatTok{0.20}\NormalTok{, }\StringTok{\textquotesingle{}beta\_m\textquotesingle{}}\NormalTok{: }\FloatTok{1.3}\NormalTok{,}
                \StringTok{\textquotesingle{}alfa\_Y\textquotesingle{}}\NormalTok{: }\FloatTok{0.15}\NormalTok{, }\StringTok{\textquotesingle{}beta\_Y\textquotesingle{}}\NormalTok{: }\FloatTok{1.1}\NormalTok{,}
                \StringTok{\textquotesingle{}k\_micro\textquotesingle{}}\NormalTok{: }\FloatTok{1.8}\NormalTok{, }\StringTok{\textquotesingle{}k\_inter\textquotesingle{}}\NormalTok{: }\FloatTok{0.3}\NormalTok{, }\StringTok{\textquotesingle{}n\_inter\textquotesingle{}}\NormalTok{: }\FloatTok{1.5}\NormalTok{,}
                \StringTok{\textquotesingle{}a\textquotesingle{}}\NormalTok{: }\FloatTok{0.75}\NormalTok{, }\StringTok{\textquotesingle{}F\_max\textquotesingle{}}\NormalTok{: }\FloatTok{0.45}\NormalTok{, }\StringTok{\textquotesingle{}F\_L\textquotesingle{}}\NormalTok{: }\FloatTok{0.25}\NormalTok{, }\StringTok{\textquotesingle{}m\_L\textquotesingle{}}\NormalTok{: }\FloatTok{0.06}\NormalTok{,}
                \StringTok{\textquotesingle{}B\_max\textquotesingle{}}\NormalTok{: }\DecValTok{180}\NormalTok{, }\StringTok{\textquotesingle{}n\_B\textquotesingle{}}\NormalTok{: }\FloatTok{0.75}
\NormalTok{            \},}
            \StringTok{\textquotesingle{}estiercol\textquotesingle{}}\NormalTok{: \{}
                \StringTok{\textquotesingle{}m\_base\textquotesingle{}}\NormalTok{: }\FloatTok{0.10}\NormalTok{, }\StringTok{\textquotesingle{}Y\_B\_base\textquotesingle{}}\NormalTok{: }\FloatTok{0.55}\NormalTok{, }\StringTok{\textquotesingle{}Y\_L\textquotesingle{}}\NormalTok{: }\FloatTok{0.42}\NormalTok{,}
                \StringTok{\textquotesingle{}alfa\_m\textquotesingle{}}\NormalTok{: }\FloatTok{0.25}\NormalTok{, }\StringTok{\textquotesingle{}beta\_m\textquotesingle{}}\NormalTok{: }\FloatTok{1.4}\NormalTok{,}
                \StringTok{\textquotesingle{}alfa\_Y\textquotesingle{}}\NormalTok{: }\FloatTok{0.20}\NormalTok{, }\StringTok{\textquotesingle{}beta\_Y\textquotesingle{}}\NormalTok{: }\FloatTok{1.2}\NormalTok{,}
                \StringTok{\textquotesingle{}k\_micro\textquotesingle{}}\NormalTok{: }\FloatTok{2.5}\NormalTok{, }\StringTok{\textquotesingle{}k\_inter\textquotesingle{}}\NormalTok{: }\FloatTok{0.4}\NormalTok{, }\StringTok{\textquotesingle{}n\_inter\textquotesingle{}}\NormalTok{: }\FloatTok{1.5}\NormalTok{,}
                \StringTok{\textquotesingle{}a\textquotesingle{}}\NormalTok{: }\FloatTok{0.7}\NormalTok{, }\StringTok{\textquotesingle{}F\_max\textquotesingle{}}\NormalTok{: }\FloatTok{0.4}\NormalTok{, }\StringTok{\textquotesingle{}F\_L\textquotesingle{}}\NormalTok{: }\FloatTok{0.2}\NormalTok{, }\StringTok{\textquotesingle{}m\_L\textquotesingle{}}\NormalTok{: }\FloatTok{0.07}\NormalTok{,}
                \StringTok{\textquotesingle{}B\_max\textquotesingle{}}\NormalTok{: }\DecValTok{160}\NormalTok{, }\StringTok{\textquotesingle{}n\_B\textquotesingle{}}\NormalTok{: }\FloatTok{0.75}
\NormalTok{            \},}
            \StringTok{\textquotesingle{}lodos\_digeridos\textquotesingle{}}\NormalTok{: \{}
                \StringTok{\textquotesingle{}m\_base\textquotesingle{}}\NormalTok{: }\FloatTok{0.12}\NormalTok{, }\StringTok{\textquotesingle{}Y\_B\_base\textquotesingle{}}\NormalTok{: }\FloatTok{0.68}\NormalTok{, }\StringTok{\textquotesingle{}Y\_L\textquotesingle{}}\NormalTok{: }\FloatTok{0.42}\NormalTok{,}
                \StringTok{\textquotesingle{}alfa\_m\textquotesingle{}}\NormalTok{: }\FloatTok{0.30}\NormalTok{, }\StringTok{\textquotesingle{}beta\_m\textquotesingle{}}\NormalTok{: }\FloatTok{1.5}\NormalTok{,}
                \StringTok{\textquotesingle{}alfa\_Y\textquotesingle{}}\NormalTok{: }\FloatTok{0.25}\NormalTok{, }\StringTok{\textquotesingle{}beta\_Y\textquotesingle{}}\NormalTok{: }\FloatTok{1.3}\NormalTok{,}
                \StringTok{\textquotesingle{}k\_micro\textquotesingle{}}\NormalTok{: }\FloatTok{1.5}\NormalTok{, }\StringTok{\textquotesingle{}k\_inter\textquotesingle{}}\NormalTok{: }\FloatTok{0.25}\NormalTok{, }\StringTok{\textquotesingle{}n\_inter\textquotesingle{}}\NormalTok{: }\FloatTok{1.5}\NormalTok{,}
                \StringTok{\textquotesingle{}a\textquotesingle{}}\NormalTok{: }\FloatTok{0.65}\NormalTok{, }\StringTok{\textquotesingle{}F\_max\textquotesingle{}}\NormalTok{: }\FloatTok{0.35}\NormalTok{, }\StringTok{\textquotesingle{}F\_L\textquotesingle{}}\NormalTok{: }\FloatTok{0.15}\NormalTok{, }\StringTok{\textquotesingle{}m\_L\textquotesingle{}}\NormalTok{: }\FloatTok{0.08}\NormalTok{,}
                \StringTok{\textquotesingle{}B\_max\textquotesingle{}}\NormalTok{: }\DecValTok{140}\NormalTok{, }\StringTok{\textquotesingle{}n\_B\textquotesingle{}}\NormalTok{: }\FloatTok{0.75}
\NormalTok{            \}}
\NormalTok{        \}}
        
        \CommentTok{\# Parámetros ambientales (comunes para todos los sustratos)}
\NormalTok{        parametros\_comunes }\OperatorTok{=}\NormalTok{ \{}
            \StringTok{\textquotesingle{}E\_a\textquotesingle{}}\NormalTok{: }\DecValTok{50000}\NormalTok{, }\StringTok{\textquotesingle{}E\_h\textquotesingle{}}\NormalTok{: }\DecValTok{100000}\NormalTok{, }\StringTok{\textquotesingle{}E\_l\textquotesingle{}}\NormalTok{: }\DecValTok{100000}\NormalTok{, }
            \StringTok{\textquotesingle{}T\_h\textquotesingle{}}\NormalTok{: }\DecValTok{313}\NormalTok{, }\StringTok{\textquotesingle{}T\_l\textquotesingle{}}\NormalTok{: }\DecValTok{283}\NormalTok{, }\StringTok{\textquotesingle{}T\_ref\textquotesingle{}}\NormalTok{: }\DecValTok{300}\NormalTok{, }\StringTok{\textquotesingle{}R\textquotesingle{}}\NormalTok{: }\FloatTok{8.314}\NormalTok{,}
            \StringTok{\textquotesingle{}pH\_opt\textquotesingle{}}\NormalTok{: }\FloatTok{6.0}\NormalTok{, }\StringTok{\textquotesingle{}sigma\_pH\textquotesingle{}}\NormalTok{: }\FloatTok{1.5}\NormalTok{,}
            \StringTok{\textquotesingle{}H\_opt\textquotesingle{}}\NormalTok{: }\DecValTok{65}\NormalTok{, }\StringTok{\textquotesingle{}sigma\_H\textquotesingle{}}\NormalTok{: }\DecValTok{15}\NormalTok{,}
            \StringTok{\textquotesingle{}D\_ref\textquotesingle{}}\NormalTok{: }\DecValTok{5}\NormalTok{, }\StringTok{\textquotesingle{}gamma\textquotesingle{}}\NormalTok{: }\FloatTok{0.05}
\NormalTok{        \}}
        
        \CommentTok{\# Combinar parámetros específicos y comunes}
\NormalTok{        parametros }\OperatorTok{=}\NormalTok{ \{}\OperatorTok{**}\NormalTok{parametros\_por\_sustrato[tipo\_sustrato], }\OperatorTok{**}\NormalTok{parametros\_comunes\}}
        \ControlFlowTok{return}\NormalTok{ parametros}
    
    \KeywordTok{def}\NormalTok{ factor\_temperatura(}\VariableTok{self}\NormalTok{, T):}
        \CommentTok{\# Convertir a Kelvin si está en Celsius}
        \ControlFlowTok{if}\NormalTok{ T }\OperatorTok{\textless{}} \DecValTok{100}\NormalTok{:}
\NormalTok{            T }\OperatorTok{=}\NormalTok{ T }\OperatorTok{+} \FloatTok{273.15}
            
        \CommentTok{\# Modelo termofisiológico}
\NormalTok{        p }\OperatorTok{=} \VariableTok{self}\NormalTok{.parametros}
\NormalTok{        f\_T }\OperatorTok{=}\NormalTok{ np.exp(p[}\StringTok{\textquotesingle{}E\_a\textquotesingle{}}\NormalTok{]}\OperatorTok{/}\NormalTok{p[}\StringTok{\textquotesingle{}R\textquotesingle{}}\NormalTok{] }\OperatorTok{*}\NormalTok{ (}\DecValTok{1}\OperatorTok{/}\NormalTok{p[}\StringTok{\textquotesingle{}T\_ref\textquotesingle{}}\NormalTok{] }\OperatorTok{{-}} \DecValTok{1}\OperatorTok{/}\NormalTok{T))}
\NormalTok{        f\_T }\OperatorTok{*=}\NormalTok{ (}\DecValTok{1} \OperatorTok{{-}}\NormalTok{ np.exp(p[}\StringTok{\textquotesingle{}E\_h\textquotesingle{}}\NormalTok{]}\OperatorTok{/}\NormalTok{p[}\StringTok{\textquotesingle{}R\textquotesingle{}}\NormalTok{] }\OperatorTok{*}\NormalTok{ (}\DecValTok{1}\OperatorTok{/}\NormalTok{p[}\StringTok{\textquotesingle{}T\_h\textquotesingle{}}\NormalTok{] }\OperatorTok{{-}} \DecValTok{1}\OperatorTok{/}\NormalTok{T)))}
\NormalTok{        f\_T }\OperatorTok{*=}\NormalTok{ (}\DecValTok{1} \OperatorTok{{-}}\NormalTok{ np.exp(p[}\StringTok{\textquotesingle{}E\_l\textquotesingle{}}\NormalTok{]}\OperatorTok{/}\NormalTok{p[}\StringTok{\textquotesingle{}R\textquotesingle{}}\NormalTok{] }\OperatorTok{*}\NormalTok{ (}\DecValTok{1}\OperatorTok{/}\NormalTok{T }\OperatorTok{{-}} \DecValTok{1}\OperatorTok{/}\NormalTok{p[}\StringTok{\textquotesingle{}T\_l\textquotesingle{}}\NormalTok{])))}
        \ControlFlowTok{return} \BuiltInTok{max}\NormalTok{(}\DecValTok{0}\NormalTok{, f\_T)}
    
    \KeywordTok{def}\NormalTok{ factor\_pH(}\VariableTok{self}\NormalTok{, pH):}
\NormalTok{        p }\OperatorTok{=} \VariableTok{self}\NormalTok{.parametros}
        \ControlFlowTok{return}\NormalTok{ np.exp(}\OperatorTok{{-}}\NormalTok{((pH }\OperatorTok{{-}}\NormalTok{ p[}\StringTok{\textquotesingle{}pH\_opt\textquotesingle{}}\NormalTok{])}\OperatorTok{**}\DecValTok{2}\NormalTok{)}\OperatorTok{/}\NormalTok{(}\DecValTok{2} \OperatorTok{*}\NormalTok{ p[}\StringTok{\textquotesingle{}sigma\_pH\textquotesingle{}}\NormalTok{]}\OperatorTok{**}\DecValTok{2}\NormalTok{))}
    
    \KeywordTok{def}\NormalTok{ factor\_humedad(}\VariableTok{self}\NormalTok{, H):}
\NormalTok{        p }\OperatorTok{=} \VariableTok{self}\NormalTok{.parametros}
        \ControlFlowTok{if}\NormalTok{ H }\OperatorTok{\textless{}=}\NormalTok{ p[}\StringTok{\textquotesingle{}H\_opt\textquotesingle{}}\NormalTok{]:}
            \ControlFlowTok{return}\NormalTok{ (H}\OperatorTok{/}\NormalTok{p[}\StringTok{\textquotesingle{}H\_opt\textquotesingle{}}\NormalTok{]) }\OperatorTok{*}\NormalTok{ np.exp(}\DecValTok{1} \OperatorTok{{-}}\NormalTok{ H}\OperatorTok{/}\NormalTok{p[}\StringTok{\textquotesingle{}H\_opt\textquotesingle{}}\NormalTok{])}
        \ControlFlowTok{else}\NormalTok{:}
            \ControlFlowTok{return}\NormalTok{ np.exp(}\OperatorTok{{-}}\NormalTok{((H }\OperatorTok{{-}}\NormalTok{ p[}\StringTok{\textquotesingle{}H\_opt\textquotesingle{}}\NormalTok{])}\OperatorTok{**}\DecValTok{2}\NormalTok{)}\OperatorTok{/}\NormalTok{(}\DecValTok{2} \OperatorTok{*}\NormalTok{ p[}\StringTok{\textquotesingle{}sigma\_H\textquotesingle{}}\NormalTok{]}\OperatorTok{**}\DecValTok{2}\NormalTok{))}
    
    \KeywordTok{def}\NormalTok{ factor\_densidad(}\VariableTok{self}\NormalTok{, D):}
\NormalTok{        p }\OperatorTok{=} \VariableTok{self}\NormalTok{.parametros}
        \ControlFlowTok{if}\NormalTok{ D }\OperatorTok{\textless{}}\NormalTok{ p[}\StringTok{\textquotesingle{}D\_ref\textquotesingle{}}\NormalTok{]:}
            \ControlFlowTok{return} \FloatTok{1.0}
        \ControlFlowTok{else}\NormalTok{:}
            \ControlFlowTok{return} \BuiltInTok{max}\NormalTok{(}\FloatTok{0.1}\NormalTok{, }\DecValTok{1} \OperatorTok{{-}}\NormalTok{ p[}\StringTok{\textquotesingle{}gamma\textquotesingle{}}\NormalTok{] }\OperatorTok{*}\NormalTok{ (D}\OperatorTok{/}\NormalTok{p[}\StringTok{\textquotesingle{}D\_ref\textquotesingle{}}\NormalTok{] }\OperatorTok{{-}} \DecValTok{1}\NormalTok{))}
    
    \KeywordTok{def}\NormalTok{ tasa\_alimentacion(}\VariableTok{self}\NormalTok{, B, S):}
\NormalTok{        p }\OperatorTok{=} \VariableTok{self}\NormalTok{.parametros}
        \CommentTok{\# Tasa de alimentación dependiente de disponibilidad de sustrato}
        \ControlFlowTok{return}\NormalTok{ p[}\StringTok{\textquotesingle{}F\_max\textquotesingle{}}\NormalTok{] }\OperatorTok{*}\NormalTok{ (S }\OperatorTok{/}\NormalTok{ (S }\OperatorTok{+} \DecValTok{10}\NormalTok{))}
    
    \KeywordTok{def}\NormalTok{ sistema\_ecuaciones(}\VariableTok{self}\NormalTok{, t, y, condiciones):}
        \CommentTok{\# Desempaquetar variables de estado}
\NormalTok{        B, L, S }\OperatorTok{=}\NormalTok{ y}
        
        \CommentTok{\# Obtener condiciones ambientales en tiempo t}
\NormalTok{        T }\OperatorTok{=}\NormalTok{ condiciones[}\StringTok{\textquotesingle{}temperatura\textquotesingle{}}\NormalTok{](t)}
\NormalTok{        pH }\OperatorTok{=}\NormalTok{ condiciones[}\StringTok{\textquotesingle{}pH\textquotesingle{}}\NormalTok{](t)}
\NormalTok{        H }\OperatorTok{=}\NormalTok{ condiciones[}\StringTok{\textquotesingle{}humedad\textquotesingle{}}\NormalTok{](t)}
\NormalTok{        D }\OperatorTok{=}\NormalTok{ B }\OperatorTok{/} \BuiltInTok{max}\NormalTok{(}\FloatTok{0.1}\NormalTok{, S)  }\CommentTok{\# Densidad}
        
        \CommentTok{\# Calcular factores ambientales}
\NormalTok{        f\_T }\OperatorTok{=} \VariableTok{self}\NormalTok{.factor\_temperatura(T)}
\NormalTok{        f\_pH }\OperatorTok{=} \VariableTok{self}\NormalTok{.factor\_pH(pH)}
\NormalTok{        f\_H }\OperatorTok{=} \VariableTok{self}\NormalTok{.factor\_humedad(H)}
\NormalTok{        f\_D }\OperatorTok{=} \VariableTok{self}\NormalTok{.factor\_densidad(D)}
        
        \CommentTok{\# Parámetros dependientes del tiempo}
\NormalTok{        p }\OperatorTok{=} \VariableTok{self}\NormalTok{.parametros}
\NormalTok{        t\_max }\OperatorTok{=}\NormalTok{ condiciones[}\StringTok{\textquotesingle{}tiempo\_total\textquotesingle{}}\NormalTok{]}
\NormalTok{        m }\OperatorTok{=}\NormalTok{ p[}\StringTok{\textquotesingle{}m\_base\textquotesingle{}}\NormalTok{] }\OperatorTok{*}\NormalTok{ (}\DecValTok{1} \OperatorTok{+}\NormalTok{ p[}\StringTok{\textquotesingle{}alfa\_m\textquotesingle{}}\NormalTok{] }\OperatorTok{*}\NormalTok{ (t}\OperatorTok{/}\NormalTok{t\_max)}\OperatorTok{**}\NormalTok{p[}\StringTok{\textquotesingle{}beta\_m\textquotesingle{}}\NormalTok{])}
\NormalTok{        Y\_B }\OperatorTok{=}\NormalTok{ p[}\StringTok{\textquotesingle{}Y\_B\_base\textquotesingle{}}\NormalTok{] }\OperatorTok{*}\NormalTok{ (}\DecValTok{1} \OperatorTok{+}\NormalTok{ p[}\StringTok{\textquotesingle{}alfa\_Y\textquotesingle{}}\NormalTok{] }\OperatorTok{*}\NormalTok{ (t}\OperatorTok{/}\NormalTok{t\_max)}\OperatorTok{**}\NormalTok{p[}\StringTok{\textquotesingle{}beta\_Y\textquotesingle{}}\NormalTok{])}
        
        \CommentTok{\# Tasa de alimentación}
\NormalTok{        F }\OperatorTok{=} \VariableTok{self}\NormalTok{.tasa\_alimentacion(B, S)}
        
        \CommentTok{\# Ecuaciones de biomasa estructural y lípidos}
\NormalTok{        dB\_dt }\OperatorTok{=}\NormalTok{ p[}\StringTok{\textquotesingle{}a\textquotesingle{}}\NormalTok{] }\OperatorTok{*}\NormalTok{ F }\OperatorTok{*}\NormalTok{ B}\OperatorTok{**}\NormalTok{(}\DecValTok{2}\OperatorTok{/}\DecValTok{3}\NormalTok{) }\OperatorTok{*}\NormalTok{ (}\DecValTok{1} \OperatorTok{{-}}\NormalTok{ (B}\OperatorTok{/}\NormalTok{p[}\StringTok{\textquotesingle{}B\_max\textquotesingle{}}\NormalTok{])}\OperatorTok{**}\NormalTok{p[}\StringTok{\textquotesingle{}n\_B\textquotesingle{}}\NormalTok{]) }\OperatorTok{*}\NormalTok{ f\_T }\OperatorTok{*}\NormalTok{ f\_pH }\OperatorTok{*}\NormalTok{ f\_H }\OperatorTok{*}\NormalTok{ f\_D}
\NormalTok{        dL\_dt }\OperatorTok{=}\NormalTok{ (p[}\StringTok{\textquotesingle{}F\_L\textquotesingle{}}\NormalTok{] }\OperatorTok{*}\NormalTok{ F }\OperatorTok{{-}}\NormalTok{ p[}\StringTok{\textquotesingle{}m\_L\textquotesingle{}}\NormalTok{] }\OperatorTok{*}\NormalTok{ L) }\OperatorTok{*}\NormalTok{ f\_T }\OperatorTok{*}\NormalTok{ f\_pH }\OperatorTok{*}\NormalTok{ f\_H }\OperatorTok{*}\NormalTok{ f\_D}
        
        \CommentTok{\# Degradación microbiana del sustrato}
\NormalTok{        r\_S\_micro }\OperatorTok{=}\NormalTok{ p[}\StringTok{\textquotesingle{}k\_micro\textquotesingle{}}\NormalTok{] }\OperatorTok{*} \FloatTok{0.05} \OperatorTok{*}\NormalTok{ S }\OperatorTok{*}\NormalTok{ f\_T }\OperatorTok{*}\NormalTok{ f\_pH }\OperatorTok{*}\NormalTok{ f\_H}
        
        \CommentTok{\# Consumo de sustrato}
\NormalTok{        dS\_dt }\OperatorTok{=} \OperatorTok{{-}}\NormalTok{(F }\OperatorTok{*}\NormalTok{ B}\OperatorTok{**}\NormalTok{(}\DecValTok{2}\OperatorTok{/}\DecValTok{3}\NormalTok{) }\OperatorTok{+}\NormalTok{ r\_S\_micro) }\OperatorTok{*}\NormalTok{ f\_T }\OperatorTok{*}\NormalTok{ f\_pH }\OperatorTok{*}\NormalTok{ f\_H}
        
        \ControlFlowTok{return}\NormalTok{ [dB\_dt, dL\_dt, dS\_dt]}
    
    \KeywordTok{def}\NormalTok{ calcular\_co2(}\VariableTok{self}\NormalTok{, t, B, L, S, r\_B, r\_L, condiciones):}
        \CommentTok{\# Obtener condiciones ambientales en tiempo t}
\NormalTok{        T }\OperatorTok{=}\NormalTok{ condiciones[}\StringTok{\textquotesingle{}temperatura\textquotesingle{}}\NormalTok{](t)}
\NormalTok{        pH }\OperatorTok{=}\NormalTok{ condiciones[}\StringTok{\textquotesingle{}pH\textquotesingle{}}\NormalTok{](t)}
\NormalTok{        H }\OperatorTok{=}\NormalTok{ condiciones[}\StringTok{\textquotesingle{}humedad\textquotesingle{}}\NormalTok{](t)}
\NormalTok{        D }\OperatorTok{=}\NormalTok{ B }\OperatorTok{/} \BuiltInTok{max}\NormalTok{(}\FloatTok{0.1}\NormalTok{, S)  }\CommentTok{\# Densidad}
        
        \CommentTok{\# Calcular factores ambientales}
\NormalTok{        f\_T }\OperatorTok{=} \VariableTok{self}\NormalTok{.factor\_temperatura(T)}
\NormalTok{        f\_pH }\OperatorTok{=} \VariableTok{self}\NormalTok{.factor\_pH(pH)}
\NormalTok{        f\_H }\OperatorTok{=} \VariableTok{self}\NormalTok{.factor\_humedad(H)}
\NormalTok{        f\_D }\OperatorTok{=} \VariableTok{self}\NormalTok{.factor\_densidad(D)}
        
        \CommentTok{\# Parámetros dependientes del tiempo}
\NormalTok{        p }\OperatorTok{=} \VariableTok{self}\NormalTok{.parametros}
\NormalTok{        t\_max }\OperatorTok{=}\NormalTok{ condiciones[}\StringTok{\textquotesingle{}tiempo\_total\textquotesingle{}}\NormalTok{]}
\NormalTok{        m }\OperatorTok{=}\NormalTok{ p[}\StringTok{\textquotesingle{}m\_base\textquotesingle{}}\NormalTok{] }\OperatorTok{*}\NormalTok{ (}\DecValTok{1} \OperatorTok{+}\NormalTok{ p[}\StringTok{\textquotesingle{}alfa\_m\textquotesingle{}}\NormalTok{] }\OperatorTok{*}\NormalTok{ (t}\OperatorTok{/}\NormalTok{t\_max)}\OperatorTok{**}\NormalTok{p[}\StringTok{\textquotesingle{}beta\_m\textquotesingle{}}\NormalTok{])}
\NormalTok{        Y\_B }\OperatorTok{=}\NormalTok{ p[}\StringTok{\textquotesingle{}Y\_B\_base\textquotesingle{}}\NormalTok{] }\OperatorTok{*}\NormalTok{ (}\DecValTok{1} \OperatorTok{+}\NormalTok{ p[}\StringTok{\textquotesingle{}alfa\_Y\textquotesingle{}}\NormalTok{] }\OperatorTok{*}\NormalTok{ (t}\OperatorTok{/}\NormalTok{t\_max)}\OperatorTok{**}\NormalTok{p[}\StringTok{\textquotesingle{}beta\_Y\textquotesingle{}}\NormalTok{])}
        
        \CommentTok{\# CO2 larvario}
\NormalTok{        r\_CO2\_m }\OperatorTok{=}\NormalTok{ m }\OperatorTok{*}\NormalTok{ B}
\NormalTok{        r\_CO2\_B }\OperatorTok{=}\NormalTok{ r\_B }\OperatorTok{*}\NormalTok{ Y\_B}
\NormalTok{        r\_CO2\_L }\OperatorTok{=}\NormalTok{ r\_L }\OperatorTok{*}\NormalTok{ p[}\StringTok{\textquotesingle{}Y\_L\textquotesingle{}}\NormalTok{]}
\NormalTok{        r\_CO2\_larva }\OperatorTok{=}\NormalTok{ (r\_CO2\_m }\OperatorTok{+}\NormalTok{ r\_CO2\_B }\OperatorTok{+}\NormalTok{ r\_CO2\_L) }\OperatorTok{*}\NormalTok{ f\_T }\OperatorTok{*}\NormalTok{ f\_pH }\OperatorTok{*}\NormalTok{ f\_H }\OperatorTok{*}\NormalTok{ f\_D}
        
        \CommentTok{\# CO2 microbiano}
\NormalTok{        r\_CO2\_micro }\OperatorTok{=}\NormalTok{ p[}\StringTok{\textquotesingle{}k\_micro\textquotesingle{}}\NormalTok{] }\OperatorTok{*}\NormalTok{ S }\OperatorTok{*}\NormalTok{ f\_T }\OperatorTok{*}\NormalTok{ f\_pH }\OperatorTok{*}\NormalTok{ f\_H}
        
        \CommentTok{\# Factor de interacción}
        \ControlFlowTok{if}\NormalTok{ B }\OperatorTok{\textgreater{}} \DecValTok{0} \KeywordTok{and}\NormalTok{ S }\OperatorTok{\textgreater{}} \DecValTok{0}\NormalTok{:}
\NormalTok{            f\_interaccion }\OperatorTok{=} \BuiltInTok{max}\NormalTok{(}\DecValTok{0}\NormalTok{, }\DecValTok{1} \OperatorTok{{-}}\NormalTok{ p[}\StringTok{\textquotesingle{}k\_inter\textquotesingle{}}\NormalTok{] }\OperatorTok{*}\NormalTok{ (B}\OperatorTok{/}\NormalTok{S)}\OperatorTok{**}\NormalTok{p[}\StringTok{\textquotesingle{}n\_inter\textquotesingle{}}\NormalTok{])}
        \ControlFlowTok{else}\NormalTok{:}
\NormalTok{            f\_interaccion }\OperatorTok{=} \FloatTok{1.0}
            
\NormalTok{        r\_CO2\_micro }\OperatorTok{*=}\NormalTok{ f\_interaccion}
        
        \CommentTok{\# CO2 total}
\NormalTok{        r\_CO2\_total }\OperatorTok{=}\NormalTok{ r\_CO2\_larva }\OperatorTok{+}\NormalTok{ r\_CO2\_micro}
        
        \CommentTok{\# Convertir a masa CO2}
\NormalTok{        masa\_CO2 }\OperatorTok{=}\NormalTok{ r\_CO2\_total }\OperatorTok{/} \FloatTok{0.27} \OperatorTok{*} \DecValTok{1000}  \CommentTok{\# mg CO2/día}
        
        \ControlFlowTok{return}\NormalTok{ \{}
            \StringTok{\textquotesingle{}r\_CO2\_m\textquotesingle{}}\NormalTok{: r\_CO2\_m,}
            \StringTok{\textquotesingle{}r\_CO2\_B\textquotesingle{}}\NormalTok{: r\_CO2\_B,}
            \StringTok{\textquotesingle{}r\_CO2\_L\textquotesingle{}}\NormalTok{: r\_CO2\_L,}
            \StringTok{\textquotesingle{}r\_CO2\_larva\textquotesingle{}}\NormalTok{: r\_CO2\_larva,}
            \StringTok{\textquotesingle{}r\_CO2\_micro\textquotesingle{}}\NormalTok{: r\_CO2\_micro,}
            \StringTok{\textquotesingle{}r\_CO2\_total\textquotesingle{}}\NormalTok{: r\_CO2\_total,}
            \StringTok{\textquotesingle{}masa\_CO2\textquotesingle{}}\NormalTok{: masa\_CO2,}
            \StringTok{\textquotesingle{}f\_interaccion\textquotesingle{}}\NormalTok{: f\_interaccion}
\NormalTok{        \}}
    
    \KeywordTok{def}\NormalTok{ simular(}\VariableTok{self}\NormalTok{, condiciones\_iniciales, condiciones\_ambientales, tiempo\_total, puntos\_tiempo}\OperatorTok{=}\DecValTok{100}\NormalTok{):}
        \CommentTok{\# Condiciones iniciales}
\NormalTok{        B0 }\OperatorTok{=}\NormalTok{ condiciones\_iniciales.get(}\StringTok{\textquotesingle{}B0\textquotesingle{}}\NormalTok{, }\FloatTok{1.0}\NormalTok{)      }\CommentTok{\# g biomasa estructural}
\NormalTok{        L0 }\OperatorTok{=}\NormalTok{ condiciones\_iniciales.get(}\StringTok{\textquotesingle{}L0\textquotesingle{}}\NormalTok{, }\FloatTok{0.2}\NormalTok{)      }\CommentTok{\# g lípidos}
\NormalTok{        S0 }\OperatorTok{=}\NormalTok{ condiciones\_iniciales.get(}\StringTok{\textquotesingle{}S0\textquotesingle{}}\NormalTok{, }\FloatTok{100.0}\NormalTok{)    }\CommentTok{\# g sustrato}
        
        \CommentTok{\# Funciones de condiciones ambientales}
        \ControlFlowTok{if} \BuiltInTok{callable}\NormalTok{(condiciones\_ambientales.get(}\StringTok{\textquotesingle{}temperatura\textquotesingle{}}\NormalTok{)):}
\NormalTok{            temp\_func }\OperatorTok{=}\NormalTok{ condiciones\_ambientales[}\StringTok{\textquotesingle{}temperatura\textquotesingle{}}\NormalTok{]}
        \ControlFlowTok{else}\NormalTok{:}
\NormalTok{            temp\_valor }\OperatorTok{=}\NormalTok{ condiciones\_ambientales.get(}\StringTok{\textquotesingle{}temperatura\textquotesingle{}}\NormalTok{, }\FloatTok{27.0}\NormalTok{)}
\NormalTok{            temp\_func }\OperatorTok{=} \KeywordTok{lambda}\NormalTok{ t: temp\_valor}
            
        \ControlFlowTok{if} \BuiltInTok{callable}\NormalTok{(condiciones\_ambientales.get(}\StringTok{\textquotesingle{}pH\textquotesingle{}}\NormalTok{)):}
\NormalTok{            pH\_func }\OperatorTok{=}\NormalTok{ condiciones\_ambientales[}\StringTok{\textquotesingle{}pH\textquotesingle{}}\NormalTok{]}
        \ControlFlowTok{else}\NormalTok{:}
\NormalTok{            pH\_valor }\OperatorTok{=}\NormalTok{ condiciones\_ambientales.get(}\StringTok{\textquotesingle{}pH\textquotesingle{}}\NormalTok{, }\FloatTok{6.0}\NormalTok{)}
\NormalTok{            pH\_func }\OperatorTok{=} \KeywordTok{lambda}\NormalTok{ t: pH\_valor}
            
        \ControlFlowTok{if} \BuiltInTok{callable}\NormalTok{(condiciones\_ambientales.get(}\StringTok{\textquotesingle{}humedad\textquotesingle{}}\NormalTok{)):}
\NormalTok{            humedad\_func }\OperatorTok{=}\NormalTok{ condiciones\_ambientales[}\StringTok{\textquotesingle{}humedad\textquotesingle{}}\NormalTok{]}
        \ControlFlowTok{else}\NormalTok{:}
\NormalTok{            humedad\_valor }\OperatorTok{=}\NormalTok{ condiciones\_ambientales.get(}\StringTok{\textquotesingle{}humedad\textquotesingle{}}\NormalTok{, }\FloatTok{65.0}\NormalTok{)}
\NormalTok{            humedad\_func }\OperatorTok{=} \KeywordTok{lambda}\NormalTok{ t: humedad\_valor}
        
        \CommentTok{\# Configurar condiciones}
\NormalTok{        condiciones }\OperatorTok{=}\NormalTok{ \{}
            \StringTok{\textquotesingle{}temperatura\textquotesingle{}}\NormalTok{: temp\_func,}
            \StringTok{\textquotesingle{}pH\textquotesingle{}}\NormalTok{: pH\_func,}
            \StringTok{\textquotesingle{}humedad\textquotesingle{}}\NormalTok{: humedad\_func,}
            \StringTok{\textquotesingle{}tiempo\_total\textquotesingle{}}\NormalTok{: tiempo\_total}
\NormalTok{        \}}
        
        \CommentTok{\# Resolver sistema de ecuaciones diferenciales}
\NormalTok{        tiempos }\OperatorTok{=}\NormalTok{ np.linspace(}\DecValTok{0}\NormalTok{, tiempo\_total, puntos\_tiempo)}
\NormalTok{        sol }\OperatorTok{=}\NormalTok{ solve\_ivp(}
            \KeywordTok{lambda}\NormalTok{ t, y: }\VariableTok{self}\NormalTok{.sistema\_ecuaciones(t, y, condiciones),}
\NormalTok{            [}\DecValTok{0}\NormalTok{, tiempo\_total],}
\NormalTok{            [B0, L0, S0],}
\NormalTok{            t\_eval}\OperatorTok{=}\NormalTok{tiempos,}
\NormalTok{            method}\OperatorTok{=}\StringTok{\textquotesingle{}RK45\textquotesingle{}}
\NormalTok{        )}
        
        \CommentTok{\# Extraer resultados}
\NormalTok{        B }\OperatorTok{=}\NormalTok{ sol.y[}\DecValTok{0}\NormalTok{]}
\NormalTok{        L }\OperatorTok{=}\NormalTok{ sol.y[}\DecValTok{1}\NormalTok{]}
\NormalTok{        S }\OperatorTok{=}\NormalTok{ sol.y[}\DecValTok{2}\NormalTok{]}
\NormalTok{        t }\OperatorTok{=}\NormalTok{ sol.t}
        
        \CommentTok{\# Calcular tasas de crecimiento}
\NormalTok{        r\_B }\OperatorTok{=}\NormalTok{ np.zeros\_like(t)}
\NormalTok{        r\_L }\OperatorTok{=}\NormalTok{ np.zeros\_like(t)}
        \ControlFlowTok{for}\NormalTok{ i }\KeywordTok{in} \BuiltInTok{range}\NormalTok{(}\BuiltInTok{len}\NormalTok{(t)}\OperatorTok{{-}}\DecValTok{1}\NormalTok{):}
\NormalTok{            r\_B[i] }\OperatorTok{=}\NormalTok{ (B[i}\OperatorTok{+}\DecValTok{1}\NormalTok{] }\OperatorTok{{-}}\NormalTok{ B[i]) }\OperatorTok{/}\NormalTok{ (t[i}\OperatorTok{+}\DecValTok{1}\NormalTok{] }\OperatorTok{{-}}\NormalTok{ t[i])}
\NormalTok{            r\_L[i] }\OperatorTok{=}\NormalTok{ (L[i}\OperatorTok{+}\DecValTok{1}\NormalTok{] }\OperatorTok{{-}}\NormalTok{ L[i]) }\OperatorTok{/}\NormalTok{ (t[i}\OperatorTok{+}\DecValTok{1}\NormalTok{] }\OperatorTok{{-}}\NormalTok{ t[i])}
\NormalTok{        r\_B[}\OperatorTok{{-}}\DecValTok{1}\NormalTok{] }\OperatorTok{=}\NormalTok{ r\_B[}\OperatorTok{{-}}\DecValTok{2}\NormalTok{]}
\NormalTok{        r\_L[}\OperatorTok{{-}}\DecValTok{1}\NormalTok{] }\OperatorTok{=}\NormalTok{ r\_L[}\OperatorTok{{-}}\DecValTok{2}\NormalTok{]}
        
        \CommentTok{\# Calcular CO2}
\NormalTok{        co2\_resultados }\OperatorTok{=}\NormalTok{ []}
        \ControlFlowTok{for}\NormalTok{ i }\KeywordTok{in} \BuiltInTok{range}\NormalTok{(}\BuiltInTok{len}\NormalTok{(t)):}
\NormalTok{            co2 }\OperatorTok{=} \VariableTok{self}\NormalTok{.calcular\_co2(t[i], B[i], L[i], S[i], r\_B[i], r\_L[i], condiciones)}
\NormalTok{            co2\_resultados.append(co2)}
        
        \CommentTok{\# Extraer componentes de CO2}
\NormalTok{        r\_CO2\_larva }\OperatorTok{=}\NormalTok{ np.array([res[}\StringTok{\textquotesingle{}r\_CO2\_larva\textquotesingle{}}\NormalTok{] }\ControlFlowTok{for}\NormalTok{ res }\KeywordTok{in}\NormalTok{ co2\_resultados])}
\NormalTok{        r\_CO2\_micro }\OperatorTok{=}\NormalTok{ np.array([res[}\StringTok{\textquotesingle{}r\_CO2\_micro\textquotesingle{}}\NormalTok{] }\ControlFlowTok{for}\NormalTok{ res }\KeywordTok{in}\NormalTok{ co2\_resultados])}
\NormalTok{        r\_CO2\_total }\OperatorTok{=}\NormalTok{ np.array([res[}\StringTok{\textquotesingle{}r\_CO2\_total\textquotesingle{}}\NormalTok{] }\ControlFlowTok{for}\NormalTok{ res }\KeywordTok{in}\NormalTok{ co2\_resultados])}
\NormalTok{        masa\_CO2 }\OperatorTok{=}\NormalTok{ np.array([res[}\StringTok{\textquotesingle{}masa\_CO2\textquotesingle{}}\NormalTok{] }\ControlFlowTok{for}\NormalTok{ res }\KeywordTok{in}\NormalTok{ co2\_resultados])}
        
        \CommentTok{\# CO2 acumulado (kg CO2 / kg biomasa final seca)}
\NormalTok{        co2\_acumulado }\OperatorTok{=}\NormalTok{ np.trapz(r\_CO2\_total, t) }\OperatorTok{/} \FloatTok{0.27}
\NormalTok{        biomasa\_final\_seca }\OperatorTok{=}\NormalTok{ B[}\OperatorTok{{-}}\DecValTok{1}\NormalTok{] }\OperatorTok{+}\NormalTok{ L[}\OperatorTok{{-}}\DecValTok{1}\NormalTok{]}
\NormalTok{        co2\_por\_biomasa }\OperatorTok{=}\NormalTok{ co2\_acumulado }\OperatorTok{/}\NormalTok{ biomasa\_final\_seca}
        
\NormalTok{        resultados }\OperatorTok{=}\NormalTok{ \{}
            \StringTok{\textquotesingle{}tiempo\textquotesingle{}}\NormalTok{: t,}
            \StringTok{\textquotesingle{}biomasa\textquotesingle{}}\NormalTok{: B,}
            \StringTok{\textquotesingle{}lipidos\textquotesingle{}}\NormalTok{: L,}
            \StringTok{\textquotesingle{}sustrato\textquotesingle{}}\NormalTok{: S,}
            \StringTok{\textquotesingle{}r\_B\textquotesingle{}}\NormalTok{: r\_B,}
            \StringTok{\textquotesingle{}r\_L\textquotesingle{}}\NormalTok{: r\_L,}
            \StringTok{\textquotesingle{}r\_CO2\_larva\textquotesingle{}}\NormalTok{: r\_CO2\_larva,}
            \StringTok{\textquotesingle{}r\_CO2\_micro\textquotesingle{}}\NormalTok{: r\_CO2\_micro,}
            \StringTok{\textquotesingle{}r\_CO2\_total\textquotesingle{}}\NormalTok{: r\_CO2\_total,}
            \StringTok{\textquotesingle{}masa\_CO2\textquotesingle{}}\NormalTok{: masa\_CO2,}
            \StringTok{\textquotesingle{}co2\_acumulado\textquotesingle{}}\NormalTok{: co2\_acumulado,}
            \StringTok{\textquotesingle{}co2\_por\_biomasa\textquotesingle{}}\NormalTok{: co2\_por\_biomasa,}
            \StringTok{\textquotesingle{}co2\_detallado\textquotesingle{}}\NormalTok{: co2\_resultados}
\NormalTok{        \}}
        
        \ControlFlowTok{return}\NormalTok{ resultados}
    
    \KeywordTok{def}\NormalTok{ visualizar\_resultados(}\VariableTok{self}\NormalTok{, resultados):}
        \CommentTok{"""Genera gráficos para visualizar los resultados de la simulación."""}
\NormalTok{        t }\OperatorTok{=}\NormalTok{ resultados[}\StringTok{\textquotesingle{}tiempo\textquotesingle{}}\NormalTok{]}
        
        \CommentTok{\# Crear figura con subplots}
\NormalTok{        fig, axs }\OperatorTok{=}\NormalTok{ plt.subplots(}\DecValTok{3}\NormalTok{, }\DecValTok{1}\NormalTok{, figsize}\OperatorTok{=}\NormalTok{(}\DecValTok{10}\NormalTok{, }\DecValTok{12}\NormalTok{))}
        
        \CommentTok{\# Gráfico 1: Biomasa, lípidos y sustrato}
\NormalTok{        axs[}\DecValTok{0}\NormalTok{].plot(t, resultados[}\StringTok{\textquotesingle{}biomasa\textquotesingle{}}\NormalTok{], }\StringTok{\textquotesingle{}b{-}\textquotesingle{}}\NormalTok{, label}\OperatorTok{=}\StringTok{\textquotesingle{}Biomasa estructural\textquotesingle{}}\NormalTok{)}
\NormalTok{        axs[}\DecValTok{0}\NormalTok{].plot(t, resultados[}\StringTok{\textquotesingle{}lipidos\textquotesingle{}}\NormalTok{], }\StringTok{\textquotesingle{}g{-}\textquotesingle{}}\NormalTok{, label}\OperatorTok{=}\StringTok{\textquotesingle{}Lípidos\textquotesingle{}}\NormalTok{)}
\NormalTok{        axs[}\DecValTok{0}\NormalTok{].plot(t, resultados[}\StringTok{\textquotesingle{}sustrato\textquotesingle{}}\NormalTok{], }\StringTok{\textquotesingle{}r{-}\textquotesingle{}}\NormalTok{, label}\OperatorTok{=}\StringTok{\textquotesingle{}Sustrato\textquotesingle{}}\NormalTok{)}
\NormalTok{        axs[}\DecValTok{0}\NormalTok{].set\_xlabel(}\StringTok{\textquotesingle{}Tiempo (días)\textquotesingle{}}\NormalTok{)}
\NormalTok{        axs[}\DecValTok{0}\NormalTok{].set\_ylabel(}\StringTok{\textquotesingle{}Masa (g)\textquotesingle{}}\NormalTok{)}
\NormalTok{        axs[}\DecValTok{0}\NormalTok{].set\_title(}\StringTok{\textquotesingle{}Dinámica de crecimiento y consumo de sustrato\textquotesingle{}}\NormalTok{)}
\NormalTok{        axs[}\DecValTok{0}\NormalTok{].legend()}
\NormalTok{        axs[}\DecValTok{0}\NormalTok{].grid(}\VariableTok{True}\NormalTok{)}
        
        \CommentTok{\# Gráfico 2: Tasas de CO2}
\NormalTok{        axs[}\DecValTok{1}\NormalTok{].plot(t, resultados[}\StringTok{\textquotesingle{}r\_CO2\_larva\textquotesingle{}}\NormalTok{], }\StringTok{\textquotesingle{}b{-}\textquotesingle{}}\NormalTok{, label}\OperatorTok{=}\StringTok{\textquotesingle{}CO$_2$ larvario\textquotesingle{}}\NormalTok{)}
\NormalTok{        axs[}\DecValTok{1}\NormalTok{].plot(t, resultados[}\StringTok{\textquotesingle{}r\_CO2\_micro\textquotesingle{}}\NormalTok{], }\StringTok{\textquotesingle{}g{-}\textquotesingle{}}\NormalTok{, label}\OperatorTok{=}\StringTok{\textquotesingle{}CO$_2$ microbiano\textquotesingle{}}\NormalTok{)}
\NormalTok{        axs[}\DecValTok{1}\NormalTok{].plot(t, resultados[}\StringTok{\textquotesingle{}r\_CO2\_total\textquotesingle{}}\NormalTok{], }\StringTok{\textquotesingle{}r{-}\textquotesingle{}}\NormalTok{, label}\OperatorTok{=}\StringTok{\textquotesingle{}CO$_2$ total\textquotesingle{}}\NormalTok{)}
\NormalTok{        axs[}\DecValTok{1}\NormalTok{].set\_xlabel(}\StringTok{\textquotesingle{}Tiempo (días)\textquotesingle{}}\NormalTok{)}
\NormalTok{        axs[}\DecValTok{1}\NormalTok{].set\_ylabel(}\StringTok{\textquotesingle{}Tasa de CO$_2$ (gC{-}eq/día)\textquotesingle{}}\NormalTok{)}
\NormalTok{        axs[}\DecValTok{1}\NormalTok{].set\_title(}\StringTok{\textquotesingle{}Tasas de producción de CO$_2$\textquotesingle{}}\NormalTok{)}
\NormalTok{        axs[}\DecValTok{1}\NormalTok{].legend()}
\NormalTok{        axs[}\DecValTok{1}\NormalTok{].grid(}\VariableTok{True}\NormalTok{)}
        
        \CommentTok{\# Gráfico 3: Masa de CO2}
\NormalTok{        axs[}\DecValTok{2}\NormalTok{].plot(t, resultados[}\StringTok{\textquotesingle{}masa\_CO2\textquotesingle{}}\NormalTok{], }\StringTok{\textquotesingle{}r{-}\textquotesingle{}}\NormalTok{, label}\OperatorTok{=}\StringTok{\textquotesingle{}CO$_2$ (mg/día)\textquotesingle{}}\NormalTok{)}
\NormalTok{        axs[}\DecValTok{2}\NormalTok{].set\_xlabel(}\StringTok{\textquotesingle{}Tiempo (días)\textquotesingle{}}\NormalTok{)}
\NormalTok{        axs[}\DecValTok{2}\NormalTok{].set\_ylabel(}\StringTok{\textquotesingle{}mg CO$_2$/día\textquotesingle{}}\NormalTok{)}
\NormalTok{        axs[}\DecValTok{2}\NormalTok{].set\_title(}\SpecialStringTok{f\textquotesingle{}Producción de CO$_2$ (Total: }\SpecialCharTok{\{}\NormalTok{resultados[}\StringTok{"co2\_por\_biomasa"}\NormalTok{]}\SpecialCharTok{:.2f\}}\SpecialStringTok{ kg CO$_2$/kg biomasa)\textquotesingle{}}\NormalTok{)}
\NormalTok{        axs[}\DecValTok{2}\NormalTok{].legend()}
\NormalTok{        axs[}\DecValTok{2}\NormalTok{].grid(}\VariableTok{True}\NormalTok{)}
        
\NormalTok{        plt.tight\_layout()}
        \ControlFlowTok{return}\NormalTok{ fig}
\end{Highlighting}
\end{Shaded}

\subsubsection{5.3 Requerimientos para
Implementación}\label{requerimientos-para-implementaciuxf3n}

Para implementar el modelo MBCO2-H en un entorno práctico, se requiere:

\begin{enumerate}
\def\labelenumi{\arabic{enumi}.}
\tightlist
\item
  \textbf{Software}:

  \begin{itemize}
  \tightlist
  \item
    Python 3.7+ con bibliotecas NumPy, SciPy, Matplotlib
  \item
    Alternativas: MATLAB, R o implementación en lenguajes compilados
    para mayor eficiencia
  \end{itemize}
\item
  \textbf{Hardware}:

  \begin{itemize}
  \tightlist
  \item
    Requisitos mínimos: Cualquier ordenador moderno
  \item
    Para simulaciones intensivas: CPU multi-núcleo, 8+ GB RAM
  \end{itemize}
\item
  \textbf{Datos de entrada}:

  \begin{itemize}
  \tightlist
  \item
    Caracterización del sustrato (tipo, composición)
  \item
    Condiciones ambientales (temperatura, pH, humedad)
  \item
    Condiciones iniciales (biomasa inicial, cantidad de sustrato)
  \item
    Parámetros operativos (densidad larvaria, tiempo de proceso)
  \end{itemize}
\item
  \textbf{Calibración}:

  \begin{itemize}
  \tightlist
  \item
    Se recomienda calibrar parámetros específicos mediante experimentos
    de caracterización
  \end{itemize}
\end{enumerate}

\subsubsection{5.4 Limitaciones de
Implementación}\label{limitaciones-de-implementaciuxf3n}

El modelo presenta las siguientes limitaciones prácticas:

\begin{enumerate}
\def\labelenumi{\arabic{enumi}.}
\item
  \textbf{Complejidad computacional}: El sistema de ecuaciones
  diferenciales requiere métodos numéricos que pueden ser
  computacionalmente intensivos.
\item
  \textbf{Requisitos de calibración}: Para máxima precisión, se
  recomienda calibrar parámetros específicos para cada tipo de sustrato.
\item
  \textbf{Interacciones no lineales}: Algunas interacciones complejas
  entre factores pueden no estar completamente capturadas.
\item
  \textbf{Escalabilidad}: La aplicación a escalas industriales requiere
  validación adicional.
\end{enumerate}

\subsection{6. CASOS DE APLICACIÓN}\label{casos-de-aplicaciuxf3n}

\subsubsection{6.1 Caso 1: Optimización de Condiciones
Ambientales}\label{caso-1-optimizaciuxf3n-de-condiciones-ambientales}

\textbf{Objetivo}: Determinar las condiciones óptimas de temperatura, pH
y humedad para minimizar emisiones de CO$_2$ manteniendo productividad.

\textbf{Metodología}: 1. Simulación del modelo con diferentes
combinaciones de condiciones 2. Evaluación de emisiones totales y
producción de biomasa 3. Identificación de condiciones de compromiso
óptimo

\textbf{Resultados}: - Las condiciones óptimas identificadas fueron: T =
26°C, pH = 6.2, H = 60\% - Estas condiciones reducen emisiones en 18\%
vs.~condiciones estándar - La productividad se mantiene en 95\% del
máximo

\subsubsection{6.2 Caso 2: Evaluación de Diferentes
Sustratos}\label{caso-2-evaluaciuxf3n-de-diferentes-sustratos}

\textbf{Objetivo}: Comparar emisiones de CO$_2$ para diferentes sustratos y
mezclas.

\textbf{Metodología}: 1. Simulación con parámetros específicos para
diferentes sustratos 2. Normalización de emisiones por kg de biomasa
producida 3. Evaluación de factores de emisión resultantes

\textbf{Resultados}: - Residuos alimentarios + paja (1:1): 1.45 kg
CO$_2$/kg biomasa - Alimento para pollos: 2.31 kg CO$_2$/kg biomasa -
Estiércol vacuno: 2.89 kg CO$_2$/kg biomasa - Lodos digeridos (75\%): 3.12
kg CO$_2$/kg biomasa

\subsubsection{6.3 Caso 3: Análisis de Escenarios de
Escalamiento}\label{caso-3-anuxe1lisis-de-escenarios-de-escalamiento}

\textbf{Objetivo}: Evaluar cómo las emisiones de CO$_2$ varían con el
escalamiento del proceso.

\textbf{Metodología}: 1. Simulación con diferentes escalas (1 kg, 10 kg,
100 kg, 1000 kg) 2. Incorporación de efectos de escala en ventilación y
heterogeneidad 3. Evaluación de emisiones totales y por unidad de
biomasa

\textbf{Resultados}: - Escala 1 kg: 2.3 kg CO$_2$/kg biomasa - Escala 10
kg: 2.4 kg CO$_2$/kg biomasa - Escala 100 kg: 2.7 kg CO$_2$/kg biomasa -
Escala 1000 kg: 3.1 kg CO$_2$/kg biomasa

Este incremento sugiere que los efectos de escala, principalmente
relacionados con la heterogeneidad y la transferencia de oxígeno,
afectan las emisiones de CO$_2$.

\subsection{7. CONCLUSIONES Y
RECOMENDACIONES}\label{conclusiones-y-recomendaciones}

\subsubsection{7.1 Principales
Conclusiones}\label{principales-conclusiones}

\begin{enumerate}
\def\labelenumi{\arabic{enumi}.}
\item
  \textbf{Modelo bicompartimental}: La separación explícita entre CO$_2$
  larvario y microbiano proporciona una representación más precisa del
  sistema, mejorando significativamente la capacidad predictiva en
  comparación con modelos previos.
\item
  \textbf{Factores ambientales}: La temperatura, el pH y la humedad son
  los factores con mayor impacto en las emisiones de CO$_2$, con
  sensibilidades relativas de 7.25, 14.4 y 1.72 respectivamente.
\item
  \textbf{Validación experimental}: El modelo MBCO2-H muestra una
  concordancia robusta con datos experimentales de múltiples estudios
  independientes, con errores generalmente por debajo del 5\%.
\item
  \textbf{Dinámica temporal}: El modelo captura adecuadamente la
  evolución temporal de las emisiones de CO$_2$, incluyendo las fases de
  incremento gradual, pico y posterior disminución.
\item
  \textbf{Efectos de sustrato}: La calidad del sustrato influye
  significativamente en las emisiones de CO$_2$, principalmente a través de
  los parámetros m y Y\_B, que aumentan con sustratos de menor calidad.
\end{enumerate}

\subsubsection{7.2 Avances sobre Modelos
Previos}\label{avances-sobre-modelos-previos}

El modelo MBCO2-H presenta los siguientes avances respecto a los modelos
existentes:

\begin{enumerate}
\def\labelenumi{\arabic{enumi}.}
\item
  \textbf{Diferenciación explícita} entre CO$_2$ larvario y microbiano,
  permitiendo comprender mejor las contribuciones relativas.
\item
  \textbf{Incorporación cuantitativa} de factores ambientales mediante
  funciones de respuesta fundamentadas.
\item
  \textbf{Modelado de la variabilidad temporal} de parámetros clave a lo
  largo del desarrollo larvario.
\item
  \textbf{Integración del factor de densidad} para capturar efectos de
  competencia y estrés.
\item
  \textbf{Caracterización del factor de interacción} entre larvas y
  microorganismos.
\end{enumerate}

\subsubsection{7.3 Recomendaciones
Prácticas}\label{recomendaciones-pruxe1cticas}

\begin{enumerate}
\def\labelenumi{\arabic{enumi}.}
\tightlist
\item
  \textbf{Para diseño de sistemas de bioconversión}:

  \begin{itemize}
  \tightlist
  \item
    Priorizar control preciso de temperatura (±1°C) y pH (±0.2)
  \item
    Implementar sistemas de ventilación adaptables a la tasa de emisión
  \item
    Considerar los efectos de densidad en instalaciones de gran escala
  \end{itemize}
\item
  \textbf{Para monitoreo y operación}:

  \begin{itemize}
  \tightlist
  \item
    Utilizar emisiones de CO$_2$ como indicador metabólico del sistema
  \item
    Ajustar densidades en función de la calidad del sustrato
  \item
    Mantener humedad en rango 55-65\% para balance óptimo
  \end{itemize}
\item
  \textbf{Para investigación futura}:

  \begin{itemize}
  \tightlist
  \item
    Desarrollar protocolos de calibración específicos para nuevos
    sustratos
  \item
    Validar el modelo a escalas industriales
  \item
    Incorporar efectos de variabilidad genética de las poblaciones
    larvarias
  \end{itemize}
\end{enumerate}

\subsubsection{7.4 Limitaciones y Trabajo
Futuro}\label{limitaciones-y-trabajo-futuro}

El modelo presenta las siguientes limitaciones que podrían abordarse en
trabajos futuros:

\begin{enumerate}
\def\labelenumi{\arabic{enumi}.}
\item
  \textbf{Validación a escala industrial}: Se requieren más datos
  experimentales a gran escala para validar completamente las
  predicciones de escalado.
\item
  \textbf{Interacciones microbianas específicas}: El modelo actual
  simplifica la dinámica microbiana, que podría desarrollarse con mayor
  detalle.
\item
  \textbf{Variabilidad genética}: No se consideran diferencias entre
  distintas poblaciones de H. illucens.
\item
  \textbf{Emisiones adicionales}: El enfoque actual es específico para
  CO$_2$, pero podría expandirse para incluir otros gases (NH₃, CH₄, N₂O).
\end{enumerate}

Líneas de investigación futura propuestas:

\begin{enumerate}
\def\labelenumi{\arabic{enumi}.}
\item
  \textbf{Modelado metagenómico integrado} para caracterizar con
  precisión las comunidades microbianas y su contribución al CO$_2$.
\item
  \textbf{Desarrollo de sensores y control} para implementación de
  sistemas de monitoreo continuo basados en el modelo.
\item
  \textbf{Expansión a otros gases} para evaluación de impacto ambiental
  completo.
\item
  \textbf{Integración con modelos de ciclo de vida} para análisis
  holístico de sostenibilidad.
\end{enumerate}

\subsection{8. REFERENCIAS TÉCNICAS}\label{referencias-tuxe9cnicas}

\begin{enumerate}
\def\labelenumi{\arabic{enumi}.}
\item
  Eriksen, N.T. (2022). ``Dynamic modelling of feed assimilation,
  growth, lipid accumulation, and CO2 production in black soldier fly
  larvae.'' \emph{PLOS ONE}, 10.1371/journal.pone.0276605. {[}Modelo DEB
  básico{]}
\item
  Laganaro, M., Bahrndorff, S., Eriksen, N.T. (2021). ``Growth and
  metabolic performance of black soldier fly larvae grown on low and
  high-quality substrates.'' \emph{Journal of Insects as Food and Feed}.
  {[}Parámetros experimentales{]}
\item
  Bekker, N.S. et al.~(2021). ``Impact of substrate moisture content on
  growth and metabolic performance of black soldier fly larvae.''
  {[}Efectos de humedad{]}
\item
  Boakye-Yiadom, K.A. et al.~(2022). ``Greenhouse Gas Emissions and Life
  Cycle Assessment on the Black Soldier Fly.'' \emph{Sustainability},
  10.3390/su141610456. {[}Factores de emisión{]}
\item
  Mertenat, A., Diener, S., Zurbrügg, C. (2019). ``Black Soldier Fly
  Biowaste Treatment--Assessment of Global Warming Potential.''
  \emph{Waste Management}, 10.1016/j.wasman.2018.11.040. {[}Comparación
  emisiones{]}
\item
  Frontieres Study (2021). ``Soil GHG emissions from frass
  application.'' \emph{Frontiers in Sustainable Food Systems},
  10.3389/fsufs.2021.709993. {[}Datos experimentales suelo{]}
\item
  ScienceDirect Study (2020). ``GHG reduction and C/N conversion in food
  waste by BSF.'' \emph{Environmental Pollution},
  10.1016/j.envpol.2020.114436. {[}Efectos pH{]}
\item
  Padmanabha, M. et al.~(2020). ``A comprehensive dynamic growth and
  development model of Hermetia illucens larvae.'' \emph{PLOS ONE},
  10.1371/journal.pone.0239084. {[}Funciones respiración{]}
\item
  Kok, R. (2021). ``Preliminary project design for insect production:
  part 1---overall mass and energy/heat balances.'' {[}Balances
  energéticos{]}
\item
  Parodi, A. et al.~(2020). ``Bioconversion Efficiencies, Greenhouse Gas
  and Ammonia Emissions during Black Soldier Fly Rearing--A Mass Balance
  Approach.'' \emph{Journal of Cleaner Production},
  10.1016/j.jclepro.2020.122488. {[}Balance masa gases{]}
\item
  Ratkowsky, D.A., Olley, J., McMeekin, T.A., Ball, A. (1982).
  ``Relationship between temperature and growth rate of bacterial
  cultures.'' \emph{Journal of Bacteriology}, 149(1), 1-5. {[}Modelo
  temperatura{]}
\item
  Shuler, M.L., Kargi, F. (2002). ``Bioprocess Engineering: Basic
  Concepts.'' 2nd Edition, Prentice Hall. {[}Fundamentos cinética{]}
\item
  Zhang, J., Sokhansanj, S., Wu, S., Fang, R., Yang, W., Winter, P.
  (2016). ``A comprehensive model on CO2 emission from on-farm storage
  bins.'' \emph{Drying Technology}, 34(8), 881-889. {[}Emisiones CO2
  almacenamiento{]}
\item
  Li, Q., Zheng, L., Qiu, N., Cai, H., Tomberlin, J.K., Yu, Z. (2011).
  ``Bioconversion of dairy manure by black soldier fly (Diptera:
  Stratiomyidae) for biodiesel and sugar production.'' \emph{Waste
  Management}, 31(6), 1316-1320. {[}Bioconversión{]}
\item
  Heussler, C.D., Walter, A., Oberkfell, H., Dubois, S., Oettmeier, W.,
  Spieß, A.C., Stronks, H.J., Bruno, D. (2018). ``Influence of
  substrate, temperature, and pH on the performance of a microalgal
  based feed for Hermetia illucens.'' \emph{Waste and Biomass
  Valorization}, 9, 2191-2199. {[}Efectos ambientales{]}
\end{enumerate}

\begin{center}\rule{0.5\linewidth}{0.5pt}\end{center}

\emph{Este documento técnico presenta un modelo matemático propio para
la generación de CO$_2$ en bioconversión con Hermetia illucens, integrando
y ampliando los modelos existentes para proporcionar una herramienta
predictiva más precisa y aplicable en entornos prácticos.}

\end{document}
